\documentclass[11pt,letterpaper]{article}
\usepackage[T1]{fontenc}
\usepackage[utf8]{inputenc}
\usepackage{lmodern}
\usepackage[margin=1in,headheight=15pt]{geometry}
\usepackage{amsmath,amssymb,amsthm,mathtools,mathrsfs}
\usepackage{microtype,booktabs,longtable,array,enumitem}
\usepackage[dvipsnames]{xcolor}
\usepackage{fancyhdr}
\usepackage[colorlinks=true,linkcolor=MidnightBlue,citecolor=MidnightBlue,urlcolor=MidnightBlue]{hyperref}
\usepackage[nameinlink,noabbrev]{cleveref}
\hypersetup{pdftitle={Surcomplex Trigonometry: Canonical Phases, Arbitrary-Scale Triangles, and Infinitesimal Degeneration},pdfauthor={Research exposition},pdfsubject={Surreal numbers, trigonometry, non-Archimedean geometry}}
\pagestyle{fancy}\fancyhf{}
\fancyhead[L]{\small Surcomplex trigonometry}
\fancyhead[R]{\small Phases, geometry, and degeneration}
\fancyfoot[C]{\thepage}
\setlength{\parskip}{0.3em}
\setlength{\parindent}{1.3em}
\setlist{nosep,leftmargin=2em}
\newtheorem{theorem}{Theorem}[section]
\newtheorem{lemma}[theorem]{Lemma}
\newtheorem{proposition}[theorem]{Proposition}
\newtheorem{corollary}[theorem]{Corollary}
\theoremstyle{definition}
\newtheorem{definition}[theorem]{Definition}
\newtheorem{example}[theorem]{Example}
\theoremstyle{remark}
\newtheorem{remark}[theorem]{Remark}
\newcommand{\No}{\mathbf{No}}
\newcommand{\SC}{\mathbf{SC}}
\newcommand{\Oz}{\mathbf{Oz}}
\newcommand{\R}{\mathbb R}\newcommand{\C}{\mathbb C}
\newcommand{\Z}{\mathbb Z}\newcommand{\N}{\mathbb N}
\newcommand{\OO}{\mathcal O}\newcommand{\mm}{\mathfrak m}
\newcommand{\UU}{\mathbb U}\newcommand{\HH}{\mathcal H}
\newcommand{\sh}{^{\sharp}}
\newcommand{\st}{\operatorname{st}}
\newcommand{\fin}{\operatorname{fin}}
\newcommand{\supp}{\operatorname{supp}}
\newcommand{\cis}{\operatorname{cis}}
\newcommand{\Sin}{\operatorname{Sin}}
\newcommand{\Cos}{\operatorname{Cos}}
\newcommand{\Exp}{\operatorname{Exp}}
\newcommand{\expi}{\operatorname{exp}_{\mathrm{inf}}}
\newcommand{\logi}{\operatorname{log}_{\mathrm{inf}}}
\newcommand{\Arg}{\operatorname{Arg}}
\newcommand{\arcosh}{\operatorname{arcosh}}
\newcommand{\re}{\operatorname{Re}}\newcommand{\im}{\operatorname{Im}}
\newcommand{\lc}{\operatorname{lc}}
\newcommand{\HInt}{\mathop{\int^{\!H}}\nolimits}
\newcommand{\abs}[1]{\left|#1\right|}
\newcommand{\norm}[1]{\left\lVert#1\right\rVert}
\newcolumntype{L}[1]{>{\raggedright\arraybackslash}p{#1}}
\numberwithin{equation}{section}
\begin{document}
\begin{titlepage}
\centering
\vspace*{0.8cm}
{\Large\scshape A theorem-driven development\par}
\vspace{0.7cm}
{\Huge\bfseries Surcomplex Trigonometry\par}
\vspace{0.6cm}
{\LARGE Canonical Phases,\par Arbitrary-Scale Triangles,\par and Infinitesimal Degeneration\par}
\vspace{0.8cm}
{\large September 21, 2026\par}
\vspace{0.65cm}
\begin{minipage}{0.92\textwidth}
\begin{abstract}
We develop trigonometry in the surcomplex plane $\SC=\No[i]$ while distinguishing three structures: the algebraic unit circle over the real-closed field $\No$, Taylor-lifted trigonometric functions at finite angles, and the canonical global functions associated with the omnific integers. Every surcomplex direction has a finite angle. Consequently the laws of sines and cosines, triangle reconstruction, Heron's formula, half-angle and tangent laws, the inscribed-angle theorem, and Ptolemy's theorem remain exact for triangles and circles whose lengths may be infinitesimal or infinite.

The non-Archimedean information is not discarded by taking standard parts. We prove a local valuation-isometry theorem for phase, exact valuation forms of the sine law, square-root laws for nearly flat triangles, and a scale-sensitive formula for the triangle-inequality defect. Trigonometric tangency is described by a rank-two intersection algebra whose residue pairing survives collision. A quantitative inverse-cosine theorem gives the sharp condition $v(\varepsilon)>2v(\sin\theta)$ for stable linearized inversion. Finite trigonometric polynomials, exact Fourier identities, and spherical and hyperbolic analogues complete the development.

The global normalization agrees with the established Ehrlich--Kaplan construction; it is not claimed as a new definition. We also classify phase extensions once the finite-angle map is fixed, explaining precisely why local analytic laws alone do not select that normalization. All infinite sums are justified by Hahn support conditions, rather than by nonexistent fine-topological convergence of ordinary partial sums.
\end{abstract}
\end{minipage}
\vfill
\begin{minipage}{0.92\textwidth}\small
\textbf{Status.} This article supplies explicit mathematical proofs and finite symbolic checks. It is an exposition and development, not a claim of exhaustive novelty or a formally verified proof library. The role of the three supplied manuscripts and the additional published literature is specified in \S\ref{sec:provenance}.
\end{minipage}
\end{titlepage}
\tableofcontents
\newpage

\section{The objects, the results, and their provenance}\label{sec:intro}
A surcomplex number is $z=x+iy$, with $x,y\in\No$ and $i^2=-1$. The real closedness of $\No$ gives a positive square root of every positive element and makes $\SC$ algebraically closed \cite{Conway,Gonshor}. These facts already provide a norm, rotations, circles, and polynomial intersection theory. They do not, by themselves, specify an angle parameter or the value of a trigonometric function at an infinite argument.

The central observation is that \emph{infinite distances do not require infinite geometric angles}. If $z\ne0$, its normalized direction $z/|z|$ has finite coordinates. Its ordinary direction can be corrected by an infinitesimal angle, giving an exact polar representation. This permits Euclidean trigonometry at every surreal length scale. Separately, one may ask for sine, cosine, and an exponential on the full class field. We address that question using a known canonical normalization and state the additional choice it embodies.

\subsection{Three layers that should not be confused}
We use the following distinction throughout:
\begin{enumerate}
\item \textbf{Algebraic geometry over $\No$.} Dot products, determinants, square roots, and finite polynomial identities apply without a convergence theory.
\item \textbf{Finite-angle analysis.} Ordinary analytic functions are evaluated at a real standard part plus an infinitesimal by strongly summable Taylor expansions.
\item \textbf{Global normalized trigonometry.} Purely infinite normal-form terms are removed from the angle before the finite-angle function is evaluated. This is the normalization associated with the omnific integer part.
\end{enumerate}
The first two layers suffice for every triangle theorem in the article. The third explains global periodicity, the full surcomplex exponential, and the differences between geometric angles and arbitrary surreal parameters.

\subsection{Relation to the supplied sources}\label{sec:provenance}
The supplied \emph{Surcomplex Analysis} manuscript \cite{InputAnalysis} supplies the terminology of strong summability, scalar Taylor lifts, Hahn-coherent functions, and coefficientwise contour functionals. We retain those distinctions. Its family of global exponentials is revisited in \S\ref{sec:freedom}; the member which discards purely infinite imaginary terms is identified with the published canonical normalization.

The two supplied papers entitled \emph{Finite Geometry of Hahn-Coherent Surcomplex Zeros} \cite{InputGeometry} and \emph{Finite Intersections in Surcomplex Analysis} \cite{InputIntersections} concern finite-dimensional analytic intersection algebras, conservation of multiplicity, collision-stable residues, and conditioned root lifting. Their word ``geometry'' is not a theory of metric triangles. We use their intersection-theoretic viewpoint in \S\ref{sec:collisions}, where a trigonometric tangency becomes an explicitly computed rank-two algebra. The sharp inverse-Jacobian threshold in the latter manuscript motivates the inverse-cosine estimate; we prove the particular estimate directly rather than assuming the general theorem.

The external foundational inputs are surreal normal forms and real closedness \cite{Conway,Gonshor}, Hahn summability \cite{Neumann}, and restricted analytic evaluation \cite{vdDE,Alling}. Ehrlich and Kaplan already construct canonical global trigonometric functions and a canonical surcomplex exponential using the initial integer part \cite[\S11]{EK}. Accordingly, neither their existence nor the global normalization below is presented as a new discovery. The contribution here is the connected, explicit development of geometric consequences and scale-sensitive estimates, with proofs. No assertion that these consequences have never appeared elsewhere is needed.

\subsection{Main results in usable form}
Our basic parameterization is
\begin{equation}\label{eq:mainphase}
 \OO_{\R}/2\pi\Z\ \cong\ \UU,
 \qquad \UU=\{u\in\SC:u\bar u=1\},
 \qquad [\theta]\longmapsto\cos\theta+i\sin\theta.
\end{equation}
Here $\OO_{\R}$ denotes the finite real surreals, not just the ordinary reals. For every noncollinear triangle, with the usual opposite sides and angles,
\[
 \alpha+\beta+\gamma=\pi,\qquad
 a^2=b^2+c^2-2bc\cos\alpha,\qquad
 \frac a{\sin\alpha}=\frac b{\sin\beta}=\frac c{\sin\gamma}=2R.
\]
Every quantity is exact in $\No$; in particular $R$ and the area may be infinite.

A representative genuinely non-Archimedean conclusion concerns a nearly flat triangle. Put $B=b+c$, $\delta=B-a>0$, and suppose $a$ is the largest side. Then flatness means
\[
 \delta\prec\frac{bc}{b+c},
\]
not necessarily that $\delta$ is absolutely infinitesimal. In that regime,
\begin{equation}\label{eq:flatpreview}
 \pi-\alpha\sim\sqrt{\frac{2B\delta}{bc}},\qquad
 \Delta\sim\sqrt{\frac{bcB\delta}{2}},\qquad
 R\sim\sqrt{\frac{Bbc}{8\delta}}.
\end{equation}
The symbols $\prec$ and $\sim$ are defined in the next section. Their meaning is algebraic and valuation-theoretic, not a limit along a sequence.

\section{Normal forms, summability, and analytic lifting}\label{sec:foundations}
\subsection{Size conventions and valuation}
We work in a class set theory such as NBG. Every normal-form support, formal coefficient sequence, and summation family used here is a set. A fixed calculation can be placed in a set-sized Hahn workspace; no sum indexed by all surreal numbers is used. Every integer $n$, polynomial degree, Fourier grid size, or polygon size in this article is an \emph{ordinary finite integer} unless explicitly stated otherwise.

Write a normal form with increasing exponents as
\[
 z=\sum_{\gamma\in S}z_\gamma t^\gamma,
 \qquad t^\gamma:=\omega^{-\gamma},\qquad z_\gamma\in\C,
\]
where $S\subseteq\No$ is well ordered. The notation $t^\gamma$ denotes a Conway monomial, not a new use of the exponential. For $z\ne0$, put
\[
 v(z)=\min S,\qquad \lc(z)=z_{v(z)},\qquad v(0)=+\infty.
\]
Normal-form multiplication gives $v(zw)=v(z)+v(w)$ and $v(z+w)\ge\min(v(z),v(w))$. Conjugation is coefficientwise. Set
\[
 |x+iy|=\sqrt{x^2+y^2}.
\]
Then $|zw|=|z||w|$ and $v(|z|)=v(z)$, the latter following from the leading coefficient $|z_{v(z)}|$.

The finite elements and infinitesimals are
\[
 \OO_{\SC}=\{z:|z|<n\text{ for some }n\in\N_{>0}\},\qquad
 \mm_{\SC}=\{z:|z|<1/n\text{ for all }n\in\N_{>0}\}.
\]
Intersecting with $\No$ defines $\OO_{\R}$ and $\mm_{\R}$. Every finite element is uniquely $c+\eta$, with $c\in\C$ and $\eta\in\mm_{\SC}$; $c=\st(z)$ is its standard part. Equivalently $v(z)\ge0$ characterizes finiteness and $v(z)>0$ characterizes a nonzero infinitesimal. For nonzero $b$, write
\[
 a\prec b\quad\Longleftrightarrow\quad a/b\in\mm_{\SC},
 \qquad
 a\sim b\quad\Longleftrightarrow\quad a/b-1\in\mm_{\SC}.
\]
The second notation is used only with $a,b\ne0$. Thus $a\sim b$ implies equal valuation and equal leading coefficient.

\subsection{The support condition behind every infinite expansion}
\begin{definition}[Strong summability]
A set-indexed family $(z_j)$ is strongly summable if the union of its supports is well ordered and every exponent occurs in the supports of only finitely many members. Its sum is obtained by finite addition at each exponent.
\end{definition}
The Hahn--Neumann support lemma states that sums of well-ordered supports are well ordered, with finite decompositions at each exponent. If $E$ is a well-ordered set of strictly positive exponents, the monoid
\[
 E^*=\{e_1+\cdots+e_m:m\ge0,\ e_j\in E\}
\]
is well ordered and has only finitely many representations of a given exponent by nondecreasing finite words \cite{Neumann}. Consequently, for finitely many infinitesimals $\eta_1,\ldots,\eta_d$ and any formal series $F\in\C[[X_1,\ldots,X_d]]$, the evaluation $F(\eta_1,\ldots,\eta_d)$ is strongly summable. Indeed all exponents lie in the monoid generated by their positive supports, with finite contribution at each exponent. Formal addition, multiplication, and composition with zero-constant inner series are respected.

In particular,
\begin{equation}\label{eq:inflogexp}
 \expi(\eta)=\sum_{n\ge0}\frac{\eta^n}{n!},\qquad
 \logi(1+\eta)=\sum_{n\ge1}\frac{(-1)^{n+1}\eta^n}{n}
\end{equation}
are inverse maps between $\mm_{\SC}$ and $1+\mm_{\SC}$. The formal exponential and logarithm addition identities hold there. These assertions remain valid when the value group has arbitrary rank; no cofinality of the sequence $n\,v(\eta)$ is assumed.

\subsection{Taylor lifts and the meaning of differentiation}
If $f$ is an ordinary real or complex analytic function near an ordinary point $c$, define
\begin{equation}\label{eq:lift}
 f\sh(c+\eta)=\sum_{n\ge0}\frac{f^{(n)}(c)}{n!}\eta^n,
 \qquad \eta\in\mm.
\end{equation}
For several variables use the corresponding multivariable Taylor series. This is the scalar lift used in the supplied analysis paper; it is also the standard restricted-analytic evaluation on surreal numbers \cite{InputAnalysis,vdDE}. An ordinary analytic identity lifts to an identity on matching monads because the corresponding formal identity does. In particular, compatible compositions lift.

A derivative below is a derivative of a \emph{function of a surreal or surcomplex variable}, using all positive surreal tolerances. It is not a derivation on the coefficient field which differentiates the monomials $t^\gamma$. The latter would be a different operation. Taylor lifting commutes with these function derivatives. To see the local estimate directly, expand at the point of differentiation:
\[
 f\sh(z+h)=f\sh(z)+(f')\sh(z)h+h^2G(h).
\]
For sufficiently small infinitesimal $h$, the support calculation gives a fixed lower bound on the valuation of $G(h)$, hence a fixed surreal bound on $|G(h)|$. For any positive tolerance, shrinking $|h|$ makes $|hG(h)|$ smaller than that tolerance. Repeating this argument gives all higher derivatives.

\begin{lemma}[Leading-term test for a lifted germ]\label{lem:leading}
Suppose $f$ has a zero of ordinary order $m$ at $c$, with $a=f^{(m)}(c)/m!\ne0$. For every nonzero infinitesimal $\eta$,
\[
 f\sh(c+\eta)=a\eta^m(1+\xi),\qquad \xi\in\mm_{\SC}.
\]
Thus $v(f\sh(c+\eta))=m\,v(\eta)$ and its leading coefficient is that of $a\eta^m$.
\end{lemma}
\begin{proof}
Factor the ordinary germ as $f(c+X)=aX^m(1+Xg(X))$. The evaluated series $\eta g\sh(\eta)$ has strictly positive support by the support lemma. This proves the factorization and its valuation consequences.
\end{proof}
For a real analytic germ this also proves sign preservation at infinitesimal displacements, once the first nonzero term is known. For example a nonnegative ordinary germ at an interior zero has even order and positive leading coefficient, so its lift is nonnegative on that monad. Strict ordinary inequalities with a nonzero ordinary margin are preserved simply by taking standard parts. These are the two mechanisms used for inequalities below.

\section{Finite-angle trigonometry and the unit circle}\label{sec:finite}
\subsection{Definitions and exact identities}
For a finite real $\theta=r+\varepsilon$, with $r\in\R$ and $\varepsilon\in\mm_{\R}$, define
\begin{align}\label{eq:finitesin}
 \sin\theta&=\sin r\sum_{n\ge0}\frac{(-1)^n\varepsilon^{2n}}{(2n)!}
 +\cos r\sum_{n\ge0}\frac{(-1)^n\varepsilon^{2n+1}}{(2n+1)!},\\
 \cos\theta&=\cos r\sum_{n\ge0}\frac{(-1)^n\varepsilon^{2n}}{(2n)!}
 -\sin r\sum_{n\ge0}\frac{(-1)^n\varepsilon^{2n+1}}{(2n+1)!}.\nonumber
\end{align}
Define $\cis\theta=\cos\theta+i\sin\theta$. Lower-case trigonometric notation will refer to these finite-angle functions until the global extension is defined. On finite complex arguments the same notation denotes the scalar lift of ordinary complex sine and cosine.

\begin{proposition}[Finite-angle identities]\label{prop:finiteidentities}
For finite real $\theta,\phi$,
\begin{gather*}
 \cis(\theta+\phi)=\cis\theta\,\cis\phi,\qquad
 \overline{\cis\theta}=\cis(-\theta),\qquad |\cis\theta|=1,\\
 \sin(\theta+\phi)=\sin\theta\cos\phi+\cos\theta\sin\phi,\\
 \cos(\theta+\phi)=\cos\theta\cos\phi-\sin\theta\sin\phi,\qquad
 \sin^2\theta+\cos^2\theta=1.
\end{gather*}
The derivatives are $(\sin)'=\cos$, $(\cos)'=-\sin$, and $(\cis)'=i\cis$.
\end{proposition}
\begin{proof}
For $\theta=r+\varepsilon$,
\[
 \cis\theta=e^{ir}\expi(i\varepsilon).
\]
Use ordinary addition at the real standard parts and the formal exponential addition identity at the infinitesimal parts. Conjugating and multiplying proves the modulus statement; taking real and imaginary parts gives the addition formulas. Differentiation follows from \eqref{eq:lift}.
\end{proof}
It would be incorrect to interpret the ordinary Taylor series at $0$ as a strongly summable family at every nonzero real argument: infinitely many of its terms have exponent zero. Formula \eqref{eq:finitesin} first evaluates the ordinary real function and only then performs an infinitesimal Hahn expansion.

\subsection{Every direction has an angle}
\begin{theorem}[The finite phase theorem]\label{thm:phase}
The map $\cis:\OO_{\R}\to\UU$ is a surjective group homomorphism, with kernel $2\pi\Z$. Consequently \eqref{eq:mainphase} is an isomorphism of class groups. There is also a canonical decomposition
\begin{equation}\label{eq:unitdecomp}
 \UU\cong S^1(\C)\times(\mm_{\R},+),\qquad
 u\longmapsto\left(u_0,-i\logi(u/u_0)\right),\quad u_0=\st(u).
\end{equation}
\end{theorem}
\begin{proof}
If $u=x+iy\in\UU$, then $x^2+y^2=1$, so $x,y$ are finite. Its standard part $u_0$ is an ordinary unit complex number. The element $w=u/u_0$ belongs to $1+\mm_{\SC}$ and has modulus one. Since $\bar w=w^{-1}$, the logarithm identities imply
\[
 \overline{\logi w}=\logi\bar w=\logi(w^{-1})=-\logi w.
\]
Hence $\delta=-i\logi w$ is real and infinitesimal. Choose an ordinary real $\beta$ with $e^{i\beta}=u_0$. Then $\cis(\beta+\delta)=u_0\expi(i\delta)=u$.

If $\cis(r+\varepsilon)=1$, standard parts give $e^{ir}=1$, so $r\in2\pi\Z$. Then $\expi(i\varepsilon)=1$, and applying $\logi$ gives $\varepsilon=0$. This proves the kernel statement. Finally, standard parts multiply and the infinitesimal logarithm is additive on products, so \eqref{eq:unitdecomp} is a group homomorphism. Its inverse is $(u_0,\delta)\mapsto u_0\expi(i\delta)$.
\end{proof}

\begin{corollary}[Polar form at arbitrary radius]\label{cor:polar}
Every $z\in\SC^\times$ has the representation
\[
 z=\rho\cis\theta,\qquad \rho=|z|\in\No_{>0},\quad\theta\in\OO_{\R}.
\]
The radius is unique and the angle is unique modulo $2\pi\Z$. The radius may be infinitesimal, finite, or infinite.
\end{corollary}
\begin{proof}
Apply \cref{thm:phase} to $z/|z|$. Taking moduli gives uniqueness of $\rho$; the kernel gives uniqueness of the phase class.
\end{proof}

\subsection{Inverse functions, including infinite slopes}
On $[0,\pi]$, cosine is strictly decreasing from $1$ to $-1$. One direct proof is
\[
 \cos x-\cos y=2\sin\frac{x+y}{2}\sin\frac{y-x}{2}>0\qquad(0\le x<y\le\pi),
\]
since sine is positive on $(0,\pi)$: this follows from its ordinary sign in interior standard-part monads and from its leading term at the two endpoint monads. Every $q\in[-1,1]$ occurs because $q+i\sqrt{1-q^2}$ is on the upper semicircle and \cref{thm:phase} gives an angle in $[0,\pi]$. Thus $\arccos:[-1,1]\to[0,\pi]$ is well defined. Similarly $\arcsin:[-1,1]\to[-\pi/2,\pi/2]$ is the inverse of sine on that interval.

Tangent is defined wherever the finite-angle cosine is nonzero. Its addition law follows by division in \cref{prop:finiteidentities}. It gives a strictly increasing bijection
\begin{equation}\label{eq:atanrange}
 \tan:(-\pi/2,\pi/2)\longrightarrow\No,
\end{equation}
not just onto the finite surreals. Indeed, for any $x\in\No$, the unit vector
\[
 \frac{1+ix}{\sqrt{1+x^2}}
\]
has positive real part and therefore a unique angle in that interval, of tangent $x$. For $\alpha<\beta$ in the interval, the identity
\[
 \tan\beta-\tan\alpha=\frac{\sin(\beta-\alpha)}{\cos\alpha\cos\beta}>0
\]
proves strict increase. Denote the inverse of \eqref{eq:atanrange} by $\arctan$.

\begin{proposition}[Inverse expansions and derivatives]\label{prop:inverses}
For $x\in\No$, and for $q\in(-1,1)$ respectively,
\[
 (\arctan)'(x)=\frac1{1+x^2},\qquad
 (\arcsin)'(q)=\frac1{\sqrt{1-q^2}},\qquad
 (\arccos)'(q)=-\frac1{\sqrt{1-q^2}}.
\]
For positive infinite $x$,
\begin{equation}\label{eq:ataninfinite}
 \arctan x=\frac\pi2-\sum_{n\ge0}\frac{(-1)^n x^{-2n-1}}{2n+1}.
\end{equation}
\end{proposition}
\begin{proof}
Near any indicated point the derivative of the direct trigonometric map is nonzero, and formal inversion of its local Taylor series gives the displayed derivatives. The resulting series is evaluated on a sufficiently small neighborhood; no Archimedean radius is required. Alternatively, the arctangent formula follows locally from the argument of $1+ix$. For infinite positive $x$, the angle identity $\arctan x=\pi/2-\arctan(1/x)$ follows from the normalized directions and the prescribed intervals. The ordinary formal series of arctangent evaluated at the infinitesimal $1/x$ proves \eqref{eq:ataninfinite}.
\end{proof}

\begin{theorem}[Stereographic and half-angle coordinates]\label{thm:stereo}
The rational map
\begin{equation}\label{eq:cayley}
 t\longmapsto \frac{1+it}{1-it}
 =\frac{1-t^2+2it}{1+t^2}
\end{equation}
is a bijection from $\No$ to $\UU\setminus\{-1\}$. Its inverse at $u=x+iy$ is $t=y/(1+x)$. If $\theta\in(-\pi,\pi)$ represents $u$, then $t=\tan(\theta/2)$.
\end{theorem}
\begin{proof}
The denominator never vanishes, and multiplying by its conjugate gives the second expression and modulus one. A value cannot equal $-1$. Conversely, $x^2+y^2=1$ and $x\ne-1$ give
\[
 1+t^2=\frac{(1+x)^2+y^2}{(1+x)^2}=\frac2{1+x}.
\]
Substitution recovers $x$ and $y$. The half-angle statement is the double-angle identity. In particular an infinite stereographic coordinate is entirely compatible with a finite angle infinitesimally close to $\pi$ or $-\pi$.
\end{proof}

\subsection{Argument branches and their local differential}
For $z\ne0$, choose an ordinary argument $\beta$ of $\st(z/|z|)$ and define the locally unwrapped value
\begin{equation}\label{eq:unwrap}
 \arg_{\beta}z=\beta-i\logi\left(\frac{z/|z|}{\st(z/|z|)}\right).
\end{equation}
When $h/z$ is infinitesimal, the same ordinary direction is retained and
\begin{equation}\label{eq:argdifference}
 \arg_{\beta}(z+h)-\arg_{\beta}z=\im\logi(1+h/z).
\end{equation}
To verify this, divide the two unit directions: their quotient is $(1+h/z)/|1+h/z|$. Its logarithm is $\logi(1+h/z)$ minus a real number. Taking imaginary parts proves the formula. Thus the angular differential is $d\arg z=\im(dz/z)$; the radial logarithmic differential is $d\log|z|=\re(dz/z)$.

A branch constrained to the actual interval $(-\pi,\pi]$ is different from \eqref{eq:unwrap}. At a direction infinitesimally below the negative real axis, the unwrapped value may be $\pi+\varepsilon$, whereas the bounded principal value is $-\pi+\varepsilon$. The bounded branch is discontinuous at its cut, even in the fine topology. Retaining the whole standard-part monad gives the locally analytic unwrapped branch used for differentiation. This distinction also explains the global fine-local logarithm in the supplied analysis manuscript.

\section{Canonical global trigonometry and surcomplex exponentiation}\label{sec:global}
\subsection{The finite-part projection and omnific integers}
Every real surreal has a unique additive decomposition
\begin{equation}\label{eq:split}
 x=p+r+\varepsilon,\qquad p\in\Pi,\quad r\in\R,\quad\varepsilon\in\mm_{\R},
\end{equation}
where $\Pi$ consists of zero and the normal forms supported entirely at negative $t$-exponents. Its nonzero elements are purely infinite. Put $\fin(x)=r+\varepsilon$. This is an additive projection from $\No$ onto $\OO_{\R}$, with kernel $\Pi$. It is \emph{not} multiplicative:
\[
 \fin(\omega\cdot\omega^{-1})=1,
 \qquad \fin(\omega)\fin(\omega^{-1})=0.
\]
Thus reducing phases is compatible with adding angles, not with arbitrarily reducing the factors of their products.

The omnific integers have the normal-form description
\begin{equation}\label{eq:omnific}
 \Oz=\Pi\oplus\Z.
\end{equation}
They form a discrete subring: sums and products of purely infinite normal forms have no positive exponents and no constant term, and every positive element of $\Oz$ is at least $1$. They form an integer part of $\No$. To see this, write $x=p+f$ with $f$ finite and select the ordinary integer $n$ for which $n\le f<n+1$, including the appropriate side when $f$ is infinitesimally adjacent to an integer. Then $p+n\le x<p+n+1$.

Ehrlich--Kaplan's additional characterization is that $\Oz$ is the unique integer part of $\No$ which is initial in its simplicity tree \cite[Lemma 11.1]{EK}. We use that published characterization to explain the word \emph{canonical}; we do not need to reproduce its tree-theoretic proof. As $\Pi$ is closed under multiplication and division by nonzero ordinary real constants,
\begin{equation}\label{eq:periodclass}
 2\pi\Oz=\Pi+2\pi\Z,\qquad \pi\Oz=\Pi+\pi\Z.
\end{equation}

\begin{definition}[Canonical real-surreal sine and cosine]\label{def:globaltrig}
For every $x\in\No$, define
\[
 \Sin x=\sin(\fin x),\qquad \Cos x=\cos(\fin x),\qquad
 \Phi(x)=\Cos x+i\Sin x.
\]
\end{definition}
This is exactly the integer-part normalization of \cite[\S11]{EK}: reducing $x$ modulo $2\pi\Oz$ leaves the same finite representative modulo $2\pi\Z$ as reducing $\fin x$. In particular, the infinitesimal part is retained. For $p\in\Pi$ and infinitesimal $\varepsilon$,
\[
 \Sin p=0,\qquad \Cos p=1,\qquad
 \Sin(p+\varepsilon)=\varepsilon-\varepsilon^3/6+\cdots.
\]
The operation is not merely taking the standard part of an argument.

\subsection{The canonical exponential and complex trigonometric identities}
Let $\exp_{\No}:\No\to\No_{>0}$ be the usual surreal exponential, with inverse $\log_{\No}$. We use its established ordered-group isomorphism property, its agreement with the ordinary real exponential, and its infinitesimal Taylor law
\[
 \exp_{\No}(x+h)=\exp_{\No}(x)\expi(h)\qquad(h\in\mm_{\R})
\]
\cite{Gonshor,vdDE}. Define
\begin{equation}\label{eq:globalexp}
 \Exp(x+iy)=\exp_{\No}(x)\Phi(y).
\end{equation}

\begin{theorem}[Global exponential, Euler identity, and local analyticity]\label{thm:globalexp}
The map $\Exp:(\SC,+)\to(\SC^\times,\cdot)$ is onto, extends the ordinary complex exponential, and satisfies
\begin{gather}\label{eq:expkernel}
 \ker\Exp=2\pi i\Oz,
 \qquad |\Exp z|=\exp_{\No}(\re z),\\
 \Exp(z+h)=\Exp(z)\expi(h)\quad(h\in\mm_{\SC}),\qquad
 \Exp'(z)=\Exp(z).\label{eq:explocal}
\end{gather}
For every $z\in\SC$, put
\begin{equation}\label{eq:complextrig}
 \Sin z=\frac{\Exp(iz)-\Exp(-iz)}{2i},\qquad
 \Cos z=\frac{\Exp(iz)+\Exp(-iz)}2.
\end{equation}
These extend \cref{def:globaltrig}, satisfy Euler's formula $\Exp(iz)=\Cos z+i\Sin z$, the sine and cosine addition identities, and
\[
 \Sin^2z+\Cos^2z=1,\qquad \Sin'=\Cos,\qquad\Cos'=-\Sin.
\]
\end{theorem}
\begin{proof}
Both $\fin$ and the finite-angle phase map are additive-to-multiplicative in the appropriate sense, so \eqref{eq:globalexp} is a group homomorphism. Its modulus is immediate. For any $w\ne0$, choose $w=\rho\cis\theta$ with $\theta$ finite. Then $w=\Exp(\log_{\No}\rho+i\theta)$, proving surjectivity.

If $\Exp(x+iy)=1$, the modulus forces $x=0$, and \cref{thm:phase} forces $\fin(y)\in2\pi\Z$. Equations \eqref{eq:split} and \eqref{eq:periodclass} give the kernel. For $h=a+ib$ infinitesimal, $\fin(y+b)=\fin(y)+b$. The real exponential Taylor law and the finite phase law therefore give
\[
 \Exp(z+h)=\Exp(z)\expi(a)\expi(ib)=\Exp(z)\expi(h).
\]
This is a local surcomplex power series, so it proves analyticity and the derivative. Substitution into \eqref{eq:complextrig} gives the real restriction, Euler's formula, and the derivatives. Expanding products using the exponential group law proves the addition and Pythagorean identities.
\end{proof}
This direct calculation also identifies the map with the canonical surcomplex exponential of \cite[Theorem 11.1]{EK}. The local power-series argument establishes the fine-local analytic property needed here without asserting uniqueness from a differential equation.

For real surreal $x,y$, the definitions give the familiar but globally valid formulas
\begin{align}\label{eq:trighyper}
 \Sin(x+iy)&=\Sin x\cosh y+i\Cos x\sinh y,\\
 \Cos(x+iy)&=\Cos x\cosh y-i\Sin x\sinh y,\nonumber
\end{align}
where $\cosh y=(\exp_{\No}y+\exp_{\No}(-y))/2$ and $\sinh y=(\exp_{\No}y-\exp_{\No}(-y))/2$. In particular $\cosh^2y-\sinh^2y=1$. Finite geometric angles may henceforth be written in lower case, because the two definitions agree there.

\begin{corollary}[Zero classes and periods]\label{cor:globalzeros}
On the full surcomplex plane,
\[
 \{z:\Sin z=0\}=\pi\Oz,\qquad
 \{z:\Cos z=0\}=\pi/2+\pi\Oz.
\]
All these zeros are real surreals and are simple as fine-analytic zeros. The period class of each of $\Sin$ and $\Cos$, and of their pair, is $2\pi\Oz$.
\end{corollary}
\begin{proof}
The equation $\Sin z=0$ is $\Exp(2iz)=1$, so $2iz\in2\pi i\Oz$. For cosine, use $\Cos z=\Sin(z+\pi/2)$. At a sine zero, $\Cos^2z=1$, so its derivative is nonzero; the cosine case is identical. A common period is precisely an element $T$ for which $\Exp(iT)=1$, hence $T\in2\pi\Oz$. If $T$ is a period of sine alone, evaluating at $0$ and $\pi/2$ gives $\Sin T=0$ and $\Cos T=1$, so it is a common period. The cosine case is similar.
\end{proof}
Thus $2\pi$ is the least positive period, but the period class is not the cyclic group $2\pi\Z$: it also contains every purely infinite real surreal. This is compatible with all finite-angle geometric theorems, because their angle ambiguity is still only $2\pi\Z$.

\subsection{What the analytic laws do not determine}\label{sec:freedom}
\begin{theorem}[Classification of extensions of the finite phase map]\label{thm:characters}
Group homomorphisms $\Psi:(\No,+)\to\UU$ which agree with $\cis$ on $\OO_{\R}$ correspond exactly to group homomorphisms $\chi:\Pi\to\UU$, by
\begin{equation}\label{eq:character}
 \Psi(p+f)=\chi(p)\cis f,\qquad p\in\Pi,\quad f\in\OO_{\R}.
\end{equation}
Each resulting
\[
 E_\chi(x+iy)=\exp_{\No}(x)\Psi(y)
\]
is an onto exponential group homomorphism satisfying $E_\chi'=E_\chi$ and $|E_\chi(z)|=\exp_{\No}(\re z)$. The canonical choice is $\chi=1$.
\end{theorem}
\begin{proof}
The additive direct decomposition $\No=\Pi\oplus\OO_{\R}$ gives both directions of \eqref{eq:character}. Infinitesimal changes of the argument do not change its purely infinite part. Consequently $\Psi(y+b)=\Psi(y)\expi(ib)$ for infinitesimal real $b$. The proof of \eqref{eq:explocal} applies unchanged, as does the surjectivity proof using finite phases.
\end{proof}
For example, let $\ell(p)$ be the ordinary real coefficient of $\omega$ in the normal form of $p$. For any $a\in\R$, the rule $\chi_a(p)=e^{ia\ell(p)}$ gives such a character. Then
\[
 E_{\chi_0}(i\omega)=1,\qquad E_{\chi_\pi}(i\omega)=-1.
\]
This is the family in \cite{InputAnalysis}. The classification above concerns \emph{group homomorphisms with the specified finite restriction}; it is not a classification of all solutions of a differential equation. In particular it does not answer the separate robustness questions formulated in \cite[\S11]{EK} for more general initial fields and their integer parts. It does show why agreement on ordinary complex numbers and fine-local analytic laws alone cannot force the canonical normalization.

\subsection{De Moivre, Chebyshev polynomials, and roots}
For every ordinary integer $n$ and every $z\in\SC$,
\[
 (\Cos z+i\Sin z)^n=\Cos(nz)+i\Sin(nz).
\]
For $n\ge0$, define $T_0=1$, $T_1=X$, $T_{n+1}=2XT_n-T_{n-1}$, and $U_0=1$, $U_1=2X$, $U_{n+1}=2XU_n-U_{n-1}$. The addition formulas imply
\[
 \Cos(nz)=T_n(\Cos z),\qquad
 \Sin((n+1)z)=\Sin z\,U_n(\Cos z).
\]
These are polynomial identities, so there is no division problem at a sine zero.

\begin{corollary}[All roots of a nonzero surcomplex number]\label{cor:roots}
If $w=\rho\cis\theta\ne0$ and $n\ge1$ is an ordinary integer, its $n$ distinct $n$th roots are
\[
 \rho^{1/n}\cis\left(\frac{\theta+2\pi k}{n}\right),\qquad 0\le k<n.
\]
In particular the roots of unity are exactly the ordinary complex roots of unity.
\end{corollary}
\begin{proof}
De Moivre proves that each displayed value is a root. The finite-phase kernel makes them distinct. A polynomial of degree $n$ over a field has at most $n$ distinct roots. Taking $\rho=1$ and $\theta=0$ proves the last assertion.
\end{proof}

\section{Lengths, angles, and similarities at arbitrary scales}\label{sec:metric}
For vectors $u,v\in\SC$, define
\[
 \langle u,v\rangle=\re(\bar u v),\qquad
 [u,v]=\im(\bar u v).
\]
In coordinates these are $u_xv_x+u_yv_y$ and $u_xv_y-u_yv_x$. The elementary identity
\begin{equation}\label{eq:gram}
 \langle u,v\rangle^2+[u,v]^2=|u|^2|v|^2
\end{equation}
proves Cauchy--Schwarz and the triangle inequality over $\No$. Equality in Cauchy--Schwarz for nonzero vectors means linear dependence, since $[u,v]=0$. Equality in $|u+v|\le|u|+|v|$ means that $u,v$ have the same direction, or one is zero: squaring the equality gives $\langle u,v\rangle=|u||v|$, which selects the nonnegative proportionality factor.

For nonzero $u,v$, the directed angle from $u$ to $v$ is the finite class $[\theta]\in\OO_{\R}/2\pi\Z$ determined by
\begin{equation}\label{eq:angleunit}
 \cis\theta=\frac{\bar u v}{|u||v|}.
\end{equation}
The unoriented angle is the unique $\alpha\in[0,\pi]$ with
\[
 \cos\alpha=\frac{\langle u,v\rangle}{|u||v|},\qquad
 \sin\alpha=\frac{|[u,v]|}{|u||v|}.
\]
The second equation follows from \eqref{eq:gram} and the nonnegative sine on that interval. This definition is meaningful even when one vector has infinitesimal length and the other infinite length.

\begin{proposition}[Rotations and similarities]\label{prop:rotations}
Multiplication by $\cis\theta$ preserves lengths, dot products, determinants, and oriented angles. Every orientation-preserving $\No$-linear isometry of $\No^2$ is of this form. Every orientation-preserving affine similarity is
\[
 z\longmapsto a+\lambda\cis\theta\,z,\qquad a\in\SC,\quad\lambda\in\No_{>0}.
\]
An orientation-reversing one is obtained by composing with conjugation.
\end{proposition}
\begin{proof}
For a unit $q$, $\overline{qu}\,qv=\bar u v$, proving the preservation statements. An orthogonal matrix with determinant $1$ has first column $(x,y)$ of norm one and second column $(-y,x)$, so it is multiplication by $x+iy=\cis\theta$. Divide a similarity's linear part by its positive scale factor and then apply this description. A determinant $-1$ isometry becomes determinant $1$ after conjugation.
\end{proof}

\begin{example}[An infinite leg and an infinitesimal leg]\label{ex:thinright}
Take vertices $0$, $\omega$, and $i/\omega$. The legs are $\omega$ and $\omega^{-1}$, the area is exactly $1/2$, and the hypotenuse is
\[
 a=\omega\sqrt{1+\omega^{-4}}.
\]
The nonright angles are $\arctan(\omega^{-2})$ and $\pi/2-\arctan(\omega^{-2})$. Its circumradius is $a/2\sim\omega/2$. Standard parts of the normalized small angle alone would erase information essential to both the height and area.
\end{example}

\section{Triangle trigonometry over the full surreal field}\label{sec:triangles}
Let $A,B,C\in\SC$ be noncollinear, and define
\[
 a=|B-C|,\qquad b=|C-A|,\qquad c=|A-B|,\qquad
 \Delta=\tfrac12|[B-A,C-A]|>0.
\]
Let $\alpha,\beta,\gamma\in(0,\pi)$ be the interior angles at $A,B,C$. No finiteness assumption is imposed on $a,b,c$ or on $\Delta$.

\subsection{Angle sum, cosine law, and sine law}
\begin{theorem}[The fundamental triangle laws]\label{thm:trianglelaws}
The following identities hold exactly in $\No$:
\begin{gather}
 \alpha+\beta+\gamma=\pi,\label{eq:anglesum}\\
 a^2=b^2+c^2-2bc\cos\alpha,\label{eq:cosinelaw}\\
 \Delta=\tfrac12bc\sin\alpha
       =\tfrac12ca\sin\beta
       =\tfrac12ab\sin\gamma,\label{eq:arealaw}\\
 \frac a{\sin\alpha}=\frac b{\sin\beta}
 =\frac c{\sin\gamma}=\frac{abc}{2\Delta}=2R,\label{eq:sinelaw}
\end{gather}
where $R$ is the radius of the unique circle through the vertices. The cyclic versions of \eqref{eq:cosinelaw} also hold.
\end{theorem}
\begin{proof}
By a similarity and, if necessary, reflection, take $A=0$, $B=c>0$, and $C=b\cis\alpha$ with positive imaginary part. The formula for $|C-c|^2$ gives \eqref{eq:cosinelaw}, and the determinant gives the first equality of \eqref{eq:arealaw}. Cyclic relabeling gives the others.

For the angle sum, the unit phases of the other two interior angles are
\[
 \cis\beta=\frac{c-b\cis(-\alpha)}a,\qquad
 \cis\gamma=\frac{b-c\cis(-\alpha)}a.
\]
Their imaginary parts are positive, and their real parts are the defining normalized dot products. Put $u=\cis\alpha$. Multiplying all three phases gives
\[
 \frac{u(c-bu^{-1})(b-cu^{-1})}{a^2}
 =\frac{bc(u+u^{-1})-b^2-c^2}{a^2}=-1.
\]
The finite-phase kernel implies $\alpha+\beta+\gamma\in\pi+2\pi\Z$. Since the sum lies strictly between $0$ and $3\pi$, it is $\pi$.

A circumcenter $O=X+iY$ must be equidistant from $0$ and $c$, so $X=c/2$. Equidistance from $0$ and $C$ gives
\[
 O=\frac c2+i\frac{b-c\cos\alpha}{2\sin\alpha}.
\]
This solves a nonsingular pair of linear equations, so it exists uniquely. By \eqref{eq:cosinelaw},
\[
 4|O|^2\sin^2\alpha
 =c^2\sin^2\alpha+(b-c\cos\alpha)^2=a^2.
\]
The positive square roots give $R=|O|=a/(2\sin\alpha)$. Finally \eqref{eq:arealaw} gives the remaining equalities in \eqref{eq:sinelaw}.
\end{proof}
The proof uses neither compactness nor an Archimedean parallel-postulate argument. In particular, the angle sum is an equality of finite surreals, including their infinitesimal parts, not merely an equality of standard parts.

\subsection{Reconstruction and congruence}
\begin{theorem}[Side and angle reconstruction]\label{thm:reconstruction}
Three positive surreals $a,b,c$ occur as the sides of a nondegenerate triangle if and only if
\[
 a<b+c,\qquad b<c+a,\qquad c<a+b.
\]
Such a triangle is unique up to Euclidean isometry, with the two orientations interchanged by reflection. Conversely, positive finite $\alpha,\beta,\gamma$ with sum $\pi$, and an arbitrary positive scale $k\in\No$, determine a triangle with
\[
 a=k\sin\alpha,\qquad b=k\sin\beta,\qquad c=k\sin\gamma.
\]
It is unique up to similarity; specifying one positive side fixes its scale.
\end{theorem}
\begin{proof}
Necessity of the strict side inequalities follows from the triangle inequality and its equality condition. For sufficiency put $A=0$, $B=c$, and
\[
 X=\frac{b^2+c^2-a^2}{2c},\qquad Y=\sqrt{b^2-X^2}.
\]
The radicand is positive, because
\begin{equation}\label{eq:heronfactor}
 4c^2(b^2-X^2)
 =(a+b+c)(-a+b+c)(a-b+c)(a+b-c)>0.
\end{equation}
Then $C=X+iY$ has the required distances. The coordinate equations force $X$ and $Y^2$, so they also give uniqueness up to reflection.

For angle reconstruction, first note that for $\alpha,\beta>0$ with $\alpha+\beta<\pi$,
\[
 \sin\alpha+\sin\beta
 =2\sin\frac{\alpha+\beta}{2}\cos\frac{\alpha-\beta}{2}
 >2\sin\frac{\alpha+\beta}{2}\cos\frac{\alpha+\beta}{2}
 =\sin(\alpha+\beta).
\]
The strict comparison follows from cosine monotonicity on $[0,\pi]$ and $|\alpha-\beta|<\alpha+\beta$. Since $\sin\gamma=\sin(\alpha+\beta)$, the proposed sides satisfy the strict triangle inequalities, cyclically. A triangle therefore exists. The elementary identity
\[
 \sin^2\beta+\sin^2\gamma-\sin^2\alpha
 =2\sin\beta\sin\gamma\cos\alpha
\]
follows from $\alpha=\pi-\beta-\gamma$ and the addition formulas. Its cosine law shows that the constructed angle at $A$ is exactly $\alpha$, because cosine is injective on $[0,\pi]$; similarly for the others. The first part gives uniqueness at fixed $k$, and changing $k$ is a similarity.
\end{proof}

\subsection{Heron, half-angles, tangent law, and the incircle}
Put $s=(a+b+c)/2$. All of $s$, $s-a$, $s-b$, and $s-c$ are positive.
\begin{theorem}[Secondary triangle formulas]\label{thm:secondary}
One has
\begin{gather}
 \Delta^2=s(s-a)(s-b)(s-c),\label{eq:heron}\\
 \sin\frac\alpha2=\sqrt{\frac{(s-b)(s-c)}{bc}},\qquad
 \cos\frac\alpha2=\sqrt{\frac{s(s-a)}{bc}},\label{eq:halfangles}\\
 \tan\frac\alpha2=\sqrt{\frac{(s-b)(s-c)}{s(s-a)}},\qquad
 \frac{a-b}{a+b}
 =\frac{\tan((\alpha-\beta)/2)}{\tan((\alpha+\beta)/2)}.\label{eq:tangentlaw}
\end{gather}
The incircle exists, with inradius $r=\Delta/s$. In particular
\begin{equation}\label{eq:inhalf}
 \tan\frac\alpha2=\frac{r}{s-a}.
\end{equation}
The length of the internal bisector from $A$ is
\begin{equation}\label{eq:bisector}
 \ell_A=\frac{2bc}{b+c}\cos\frac\alpha2,
\end{equation}
and it divides $BC$ in the ratio $c:b$.
\end{theorem}
\begin{proof}
In the coordinates of \cref{thm:reconstruction}, $\Delta=cY/2$. Equation \eqref{eq:heronfactor} gives \eqref{eq:heron}. Substituting the cosine law into $(1-\cos\alpha)/2$ and $(1+\cos\alpha)/2$ gives the squares in \eqref{eq:halfangles}; all half-angle functions there are positive. Their quotient proves the first formula in \eqref{eq:tangentlaw}. The sine law and the sum-to-product identities give
\[
 \frac{a-b}{a+b}
 =\frac{\sin\alpha-\sin\beta}{\sin\alpha+\sin\beta}
 =\frac{\cos((\alpha+\beta)/2)\sin((\alpha-\beta)/2)}
 {\sin((\alpha+\beta)/2)\cos((\alpha-\beta)/2)},
\]
which is the tangent law. All denominators are nonzero because $0<\alpha+\beta<\pi$.

The barycentric point $I=(aA+bB+cC)/(a+b+c)$ is strictly inside the triangle. Its distance to $BC$, computed as an absolute determinant divided by $a$, is
\[
 \frac{a}{a+b+c}\frac{2\Delta}{a}=\frac\Delta s.
\]
The two other side distances agree by the same computation. The circle of this radius centered at $I$ is tangent to the three supporting lines, with the tangency points on the segments: $I$ lies on the internal bisectors, as is also seen in the following explicit bisector computation. Equations \eqref{eq:heron} and \eqref{eq:halfangles} give \eqref{eq:inhalf}.

With $A=0$, $B=c$, $C=b\cis\alpha$, the point
\[
 D=\frac{bB+cC}{b+c}=\frac{bc}{b+c}(1+\cis\alpha)
 =\frac{2bc}{b+c}\cos(\alpha/2)\cis(\alpha/2)
\]
is in the interior of $BC$. It has $BD/DC=c/b$, its direction from $A$ bisects the angle, and its modulus gives \eqref{eq:bisector}. The corresponding bisectors meet at $I$, proving the assertion used above. Projection of an internal-bisector point onto either ray lies at positive distance along that ray; applying this at each endpoint puts each contact point on its side segment.
\end{proof}

\subsection{A concurrency theorem in trigonometric form}
\begin{theorem}[Trigonometric Ceva]\label{thm:ceva}
Let $D,E,F$ lie strictly inside $BC,CA,AB$, respectively. The lines $AD,BE,CF$ are concurrent if and only if
\begin{equation}\label{eq:ceva}
 \frac{\sin\angle BAD}{\sin\angle DAC}
 \frac{\sin\angle CBE}{\sin\angle EBA}
 \frac{\sin\angle ACF}{\sin\angle FCB}=1.
\end{equation}
\end{theorem}
\begin{proof}
Taking ratios of the areas of the two triangles sharing the altitude from $A$, and also using their area formulas at $A$, gives
\[
 \frac{BD}{DC}=\frac cb\frac{\sin\angle BAD}{\sin\angle DAC}.
\]
The cyclic formulas have side factors $a/c$ and $b/a$, which cancel in their product. It remains to check the algebraic Ceva condition
\[
 (BD/DC)(CE/EA)(AF/FB)=1.
\]
For completeness, use homogeneous barycentric coordinates $(x:y:z)$ relative to $A,B,C$. The line from $A$ through $D$ imposes $z/y=BD/DC$, the line from $B$ through $E$ imposes $x/z=CE/EA$, and the line from $C$ through $F$ imposes $y/x=AF/FB$. A common point with positive coordinates exists exactly when the product is one. These are finite linear equations over $\No$, so the usual concurrency argument needs no topological hypothesis.
\end{proof}

\section{Circles, chords, Ptolemy, and coherent arc length}\label{sec:circles}
A circle of center $a\in\SC$ and radius $R\in\No_{>0}$ is exactly
\[
 \{a+R\cis\theta:\theta\in\OO_{\R}\}.
\]
This follows from polar form, so it parameterizes \emph{all} points of the circle, including directions differing by arbitrarily small infinitesimals.

\begin{theorem}[Chord and inscribed-angle theorems]\label{thm:chords}
For finite real $\alpha,\beta$,
\begin{equation}\label{eq:chord}
 |R\cis\beta-R\cis\alpha|
 =2R\left|\sin\frac{\beta-\alpha}{2}\right|.
\end{equation}
For three distinct points $z_j=a+R\cis\theta_j$, the directed angle at $z_1$ between the lines to $z_3$ and $z_2$ is
\begin{equation}\label{eq:inscribed}
 \arg\frac{z_2-z_1}{z_3-z_1}
 \equiv\frac{\theta_2-\theta_3}{2}\pmod{\pi\Z}.
\end{equation}
In particular, points in the same segment subtend the same interior angle, and a diameter subtends a right angle.
\end{theorem}
\begin{proof}
The addition formulas give the exact factorization
\begin{equation}\label{eq:chordfactor}
 \cis\beta-\cis\alpha
 =2i\cis\frac{\alpha+\beta}{2}\sin\frac{\beta-\alpha}{2}.
\end{equation}
Taking moduli proves \eqref{eq:chord}. In the ratio of two such factors, the quotient of the sine factors is a nonzero real surreal, hence has argument $0$ modulo $\pi$. The remaining phase proves \eqref{eq:inscribed}. Selecting the interior angles on the same side of the chord gives the same representative. For opposite endpoints of a diameter the right side is $\pi/2$ modulo $\pi$, so the interior angle is $\pi/2$.
\end{proof}

\begin{theorem}[Ptolemy inequality and cyclic equality]\label{thm:ptolemy}
For any four surcomplex points,
\begin{equation}\label{eq:ptolemyineq}
 |z_1-z_3|\,|z_2-z_4|
 \le |z_1-z_2|\,|z_3-z_4|+|z_1-z_4|\,|z_2-z_3|.
\end{equation}
For four distinct points in cyclic order on a circle, equality holds.
\end{theorem}
\begin{proof}
The polynomial identity
\[
 (z_1-z_3)(z_2-z_4)
 =(z_1-z_2)(z_3-z_4)+(z_1-z_4)(z_2-z_3)
\]
and the triangle inequality prove \eqref{eq:ptolemyineq}. For cyclic points choose finite angles
$\theta_1<\theta_2<\theta_3<\theta_4<\theta_1+2\pi$ and put
\[
 u=(\theta_2-\theta_1)/2,\quad
 v=(\theta_3-\theta_2)/2,\quad
 w=(\theta_4-\theta_3)/2.
\]
These are positive and $u+v+w<\pi$. All the chord half-angle sines that occur are therefore positive. Substituting \eqref{eq:chord} reduces equality to
\[
 \sin(u+v)\sin(v+w)
 =\sin u\sin w+\sin v\sin(u+v+w).
\]
Expand the two sides by the addition laws. The coefficients of $\sin u\cos w$ and $\cos u\sin w$ agree, and the remaining difference is
$\sin u\sin w(\cos^2v+\sin^2v-1)=0$. Thus equality holds, including when some consecutive arcs are infinitesimal.
\end{proof}

\begin{corollary}[Finite regular polygons at surreal scale]\label{cor:polygon}
For an ordinary integer $n\ge3$, the regular $n$-gon with vertices $a+R\cis(2\pi k/n)$ has side, perimeter, and area
\[
 2R\sin(\pi/n),\qquad
 2nR\sin(\pi/n),\qquad
 \frac n2R^2\sin(2\pi/n),
\]
respectively. These formulas hold for every positive surreal radius $R$.
\end{corollary}
\begin{proof}
The chord theorem gives the side length. Sum the finitely many equal sides and the finitely many positively oriented center triangles, each having determinant area $R^2\sin(2\pi/n)/2$.
\end{proof}

\subsection{Arc and sector functionals: the necessary convention}
There is a useful notion of arc length, but its construction must be specified. Following the supplied analysis manuscript \cite{InputAnalysis}, a uniformly supported Hahn family $F(s)=\sum_\gamma t^\gamma f_\gamma(s)$ along an ordinary parameter interval is integrated coefficientwise:
\[
 \HInt_a^b F(s)\,ds
 =\sum_\gamma t^\gamma\int_a^b f_\gamma(s)\,ds.
\]
If finite endpoints have infinitesimal displacements, use the Taylor corrections of an analytic primitive at their standard parts. Equivalently, whenever a coefficientwise analytic primitive $P$ is available, the endpoint functional is $P\sh(\beta)-P\sh(\alpha)$. This is not an assertion that a nonconstant ordinary real-parameter path is continuous in the full fine topology.

\begin{proposition}[Coherent circular arc and sector]\label{prop:arc}
Fix $a\in\SC$, $R>0$, and finite real $\alpha\le\beta$. For the unwrapped parameterization $\gamma(s)=a+R\cis s$, the coefficientwise arc-length and signed sector-area functionals are
\begin{equation}\label{eq:arclength}
 L=R(\beta-\alpha),\qquad
 A_{\mathrm{sector}}=\frac{R^2}{2}(\beta-\alpha).
\end{equation}
They count turns according to the chosen unwrapped interval. In particular one turn has length $2\pi R$ and sector area $\pi R^2$.
\end{proposition}
\begin{proof}
The pulled-back speed is $|\gamma'(s)|=|iR\cis s|=R$. The signed area integrand is
\[
 \tfrac12\im\bigl(\overline{\gamma(s)-a}\,\gamma'(s)\bigr)=R^2/2.
\]
Both are constant Hahn data, so their primitives are $Rs$ and $R^2s/2$. Evaluation at the finite endpoints proves \eqref{eq:arclength}. The same calculation justifies the infinitesimal endpoint corrections without an appeal to a Riemann-sum limit.
\end{proof}
For an infinitesimal positive central angle $\delta$, arc length and chord length are relatively equivalent:
$2R\sin(\delta/2)\sim R\delta$. The relative discrepancy is quadratic, independently of whether $R$ itself is tiny or enormous.

\section{Inequalities, angular valuations, and scale sensitivity}\label{sec:valuation}
\subsection{Small-angle inequalities without a limit shortcut}
\begin{proposition}[Classical local inequalities and exact leading orders]\label{prop:smallangles}
For finite $0<x<\pi/2$,
\begin{equation}\label{eq:smallineq}
 0<\sin x<x<\tan x.
\end{equation}
For nonzero infinitesimal $\delta$,
\begin{gather}\label{eq:smallequiv}
 \sin\delta\sim\delta,\qquad \tan\delta\sim\delta,\qquad
 1-\cos\delta\sim\delta^2/2,\\
 v(\sin\delta)=v(\tan\delta)=v(\delta),\qquad
 v(1-\cos\delta)=2v(\delta).\label{eq:smallval}
\end{gather}
\end{proposition}
\begin{proof}
For $\st(x)\in(0,\pi/2)$, all strict inequalities have positive ordinary margins and therefore lift. If $x$ is positive infinitesimal, the leading terms
\[
 x-\sin x=x^3/6+\cdots,\qquad \tan x-x=x^3/3+\cdots
\]
are positive by \cref{lem:leading}. If $x$ is infinitesimally below $\pi/2$, the comparison $\sin x<x$ has ordinary margin $\pi/2-1$, while $\tan x$ is positive infinite, by the reciprocal tangent at the small complementary angle. This covers every possible standard part. The leading terms of sine, tangent, and $1-\cos$ at zero also prove \eqref{eq:smallequiv} and \eqref{eq:smallval} directly.
\end{proof}
These asymptotic equivalences are statements about a single arbitrary infinitesimal. They are stronger than standard-part equalities and do not mean that a sequence of real arguments converges in the fine topology.

\begin{theorem}[Local valuation isometry of phase]\label{thm:phaseisometry}
If $\theta,\phi$ are finite real angles and $0\ne\theta-\phi\in\mm_{\R}$, then
\begin{equation}\label{eq:phaseisometry}
 |\cis\theta-\cis\phi|\sim|\theta-\phi|,
 \qquad v(\cis\theta-\cis\phi)=v(\theta-\phi).
\end{equation}
More precisely,
\[
 \frac{\cis\theta-\cis\phi}{i\cis\phi\,(\theta-\phi)}=1+\text{an infinitesimal}.
\]
For $z\ne0$ and a nonzero infinitesimal rotation $\delta$,
\begin{equation}\label{eq:rotdisplacement}
 |\cis\delta\,z-z|
 =2|z|\,|\sin(\delta/2)|\sim |z|\,|\delta|.
\end{equation}
\end{theorem}
\begin{proof}
Factor $\cis\theta-\cis\phi=\cis\phi\bigl(\expi(i(\theta-\phi))-1\bigr)$. The first nonzero term of the last factor is $i(\theta-\phi)$, with an infinitesimal relative remainder. Taking moduli and valuations proves the first claims. Equation \eqref{eq:rotdisplacement} is the chord theorem followed by \eqref{eq:smallequiv}.
\end{proof}
For example, a rotation by $\omega^{-1}$ of a vector of length $\omega$ moves its endpoint by a distance relatively equivalent to $1$. A rotation that is infinitesimal as an angle need not be infinitesimal as a displacement. Replacing $\omega$ by $\omega^2$ makes that displacement infinite.

\subsection{The sine law as an exact valuation identity}
For an interior angle $\alpha\in(0,\pi)$ write
\[
 d_\alpha=\min(\alpha,\pi-\alpha).
\]
It measures the distance to a degenerate angle, and belongs to $(0,\pi/2]$.

\begin{theorem}[Valued triangle trigonometry]\label{thm:valuedsine}
Every nondegenerate surcomplex triangle satisfies
\begin{gather}\label{eq:valuedsine}
 v(a)-v(d_\alpha)=v(b)-v(d_\beta)=v(c)-v(d_\gamma)=v(2R),\\
 v(\Delta)=v(b)+v(c)+v(d_\alpha),\label{eq:valuedarea}
\end{gather}
with cyclic versions of \eqref{eq:valuedarea}. The least valuation among $a,b,c$ is attained at least twice. If $\alpha$ and $\beta$ are both infinitesimal, then
\begin{equation}\label{eq:tinyanglesratio}
 \frac{\alpha}{\beta}\sim\frac ab.
\end{equation}
\end{theorem}
\begin{proof}
If $d_\alpha$ is infinitesimal, $\sin\alpha=\sin d_\alpha\sim d_\alpha$. If it is appreciable, both $\sin\alpha$ and $d_\alpha$ have positive ordinary standard part, hence valuation zero. Thus $v(\sin\alpha)=v(d_\alpha)$ in every case. Taking valuations in the sine and area laws proves \eqref{eq:valuedsine} and \eqref{eq:valuedarea}; nonzero ordinary constants have valuation zero.

Suppose, for example, $v(a)<v(b)$ and $v(a)<v(c)$. Then $(b+c)/a$ is infinitesimal, contradicting $a<b+c$. Thus there cannot be a unique least side valuation. Finally, when the two angles are infinitesimal, divide the exact sine law by the two relative equivalences $\sin\alpha\sim\alpha$ and $\sin\beta\sim\beta$ to obtain \eqref{eq:tinyanglesratio}.
\end{proof}
The conclusion about side valuations is not the assertion that all three sides have the same scale. Nor must one of $\sin\alpha,\sin\beta,\sin\gamma$ be appreciable: in a nearly flat triangle, one angle is infinitesimally close to $\pi$ and the other two are infinitesimal, so all three sines are infinitesimal.

\subsection{The triangle-inequality defect}
\begin{theorem}[A quadratic defect formula]\label{thm:defect}
Let $u,v\ne0$, $a=|u|$, $b=|v|$, and let their unoriented angle be a nonzero infinitesimal $\delta$. Set $D=a+b-|u+v|$. Then
\begin{equation}\label{eq:defectexact}
 D=\frac{4ab\sin^2(\delta/2)}{a+b+|u+v|},
\end{equation}
and
\begin{equation}\label{eq:defectasym}
 D\sim\frac{ab}{2(a+b)}\delta^2,
 \qquad v(D)=v(a)+v(b)+2v(\delta)-v(a+b).
\end{equation}
No comparability assumption on $a$ and $b$ is required.
\end{theorem}
\begin{proof}
Square and subtract:
\[
 (a+b)^2-|u+v|^2=2ab(1-\cos\delta)=4ab\sin^2(\delta/2).
\]
Division by the positive sum of the square roots proves \eqref{eq:defectexact}. Also
\[
 \frac{|u+v|^2}{(a+b)^2}
 =1-\frac{4ab}{(a+b)^2}\sin^2(\delta/2).
\]
Since $0<4ab/(a+b)^2\le1$, the subtracted term is infinitesimal. The positive square root is $1+$ an infinitesimal. Hence the denominator in \eqref{eq:defectexact} is relatively equivalent to $2(a+b)$, while $4\sin^2(\delta/2)\sim\delta^2$. This proves the relative formula and then its valuation.
\end{proof}
Unlike a bound that replaces $a+b$ by a real constant, this expression detects the effective smaller scale when $a$ and $b$ differ by many orders of magnitude.

\section{Nearly flat triangles and square-root degeneration}\label{sec:flat}
\begin{theorem}[Exact gap identity and the flatness threshold]\label{thm:flat}
Let $a$ be a largest side of a nondegenerate triangle and put
\[
 B=b+c,\qquad \delta=B-a>0,\qquad h=\pi-\alpha.
\]
Then $0<\delta\le\min(b,c)$ and
\begin{equation}\label{eq:gapidentity}
 4bc\sin^2(h/2)=\delta(2B-\delta).
\end{equation}
The angle defect $h$ is infinitesimal if and only if
\begin{equation}\label{eq:flatcriterion}
 \delta\prec bc/B.
\end{equation}
When these equivalent conditions hold, all three equivalents in \eqref{eq:flatpreview} hold, and
\begin{align}\label{eq:flatvaluations}
 v(h)&=\tfrac12\bigl(v(\delta)+v(B)-v(b)-v(c)\bigr),\\
 v(\Delta)&=\tfrac12\bigl(v(\delta)+v(B)+v(b)+v(c)\bigr),\nonumber\\
 v(R)&=\tfrac12\bigl(v(B)+v(b)+v(c)-v(\delta)\bigr).\nonumber
\end{align}
\end{theorem}
\begin{proof}
Because $a\ge\max(b,c)$, subtraction gives $\delta\le\min(b,c)$. The cosine law with $\alpha=\pi-h$ gives
\[
 B^2-a^2=2bc(1-\cos h)=4bc\sin^2(h/2).
\]
As $a=B-\delta$, this is \eqref{eq:gapidentity}. The factor $2-\delta/B$ lies between $3/2$ and $2$, since $\delta\le\min(b,c)\le B/2$. Therefore \eqref{eq:gapidentity} says that $\sin^2(h/2)$ is infinitesimal exactly when $\delta B/(bc)$ is infinitesimal. As $0<h/2<\pi/2$, its sine is infinitesimal exactly when the angle itself is. This proves the criterion.

Under the criterion, $\delta/B$ is infinitesimal: indeed
\[
 \frac\delta B=\frac{\delta B}{bc}\frac{bc}{B^2},\qquad 0<bc/B^2\le1/4.
\]
Equation \eqref{eq:gapidentity} and $\sin(h/2)\sim h/2$ give
\[
 bc\,h^2\sim2B\delta.
\]
Positive square roots yield the first equivalent in \eqref{eq:flatpreview}. The area law gives $\Delta=\frac{bc}{2}\sin h\sim\frac{bc}{2}h$, proving the second. Finally $a\sim B$ and $R=a/(2\sin h)\sim B/(2h)$ give the third. Taking valuations proves \eqref{eq:flatvaluations}.
\end{proof}
This theorem keeps track of a nonlinear geometric singularity. The square root of the side gap controls the angular defect, while its reciprocal controls the circumradius. The threshold $bc/(b+c)$ is comparable to the smaller of $b$ and $c$, so the gap must be small relative to the shorter adjacent side, not merely relative to the total length.

\begin{example}[A symmetric infinitesimal flattening]\label{ex:flat}
Let $0<\tau\in\mm_{\R}$, $b=c=1$, and $a=2-\tau$. The altitude to side $a$, area, and circumradius are exactly
\begin{equation}\label{eq:isosc}
 H=\sqrt{\tau-\tau^2/4},\qquad
 \Delta=(1-\tau/2)H,\qquad R=\frac1{2H}.
\end{equation}
The large-angle defect and the two other angles satisfy
\[
 h=2\arccos(1-\tau/2),\qquad \beta=\gamma=h/2.
\]
Their summable expansions include
\begin{align}\label{eq:isoexp}
 h&=2\sqrt\tau\left(1+\frac\tau{24}+\frac{3\tau^2}{640}
                         +\frac{5\tau^3}{7168}+\cdots\right),\\
 H&=\sqrt\tau\left(1-\frac\tau8-\frac{\tau^2}{128}
                         -\frac{\tau^3}{1024}+\cdots\right),\nonumber\\
 R&=\frac1{2\sqrt\tau}\left(1+\frac\tau8+\frac{3\tau^2}{128}
                         +\frac{5\tau^3}{1024}+\cdots\right).\nonumber
\end{align}
To derive the first, use $h=4\arcsin(\sqrt\tau/2)$ and the infinitesimal inverse-sine series. The other two are binomial expansions of \eqref{eq:isosc}. In this family the side lengths have ordinary scale, the altitude and area have square-root infinitesimal scale, and the circumradius is infinite.
\end{example}

\begin{example}[An appreciable side gap can still be a flat degeneration]
Take $b=c=\omega^2$ and $\delta=1$, so $a=2\omega^2-1$. The criterion \eqref{eq:flatcriterion} holds because $bc/B=\omega^2/2$. Thus
\[
 h\sim2/\omega,\qquad \Delta\sim\omega^3,\qquad R\sim\omega^3/2.
\]
The triangle is angularly nearly flat although its area is infinite and its gap is the ordinary number $1$. This is a reason to state degeneration hypotheses relatively rather than as absolute smallness assumptions.
\end{example}

\section{Trigonometric intersections, collisions, and sharp inversion}\label{sec:collisions}
This section makes contact with the finite-intersection viewpoint of the two supplied companion papers \cite{InputGeometry,InputIntersections}. Rather than invoke their general division and residue theorems, we compute the relevant quadratic algebra explicitly. It already displays conservation of multiplicity, a nondegenerate pairing at collision, and an inverse-Jacobian stability threshold.

\subsection{A complete linear-trigonometric intersection problem}
\begin{theorem}[Line--circle intersection and trigonometric tangency]\label{thm:linecircle}
Let $A,B,C\in\No$, with $R=\sqrt{A^2+B^2}>0$, and let $D=R^2-C^2$. The equation
\begin{equation}\label{eq:linecircle}
 A\cos\theta+B\sin\theta=C
\end{equation}
has, modulo $2\pi\Z$ in finite angles, exactly two solutions if $D>0$, exactly one if $D=0$, and none if $D<0$. The single solution when $D=0$ has multiplicity two. For $D>0$, the corresponding unit points are
\begin{equation}\label{eq:intersectionpoints}
 (X,Y)=\frac{C}{R^2}(A,B)
 \mathbin{\pm}\frac{\sqrt D}{R^2}(-B,A).
\end{equation}
At either simple solution, the absolute angular derivative of the left side of \eqref{eq:linecircle} is $\sqrt D$.
\end{theorem}
\begin{proof}
Choose a finite $\phi$ with $A+iB=R\cis\phi$. The left side is $R\cos(\theta-\phi)$. The range and strict monotonicity of cosine on a semicircle give the stated solution counts. At $C/R=\pm1$, the first derivative vanishes and the second derivative is $\mp R\ne0$, proving multiplicity two.

For a direct algebraic calculation, use orthonormal coordinates
\[
 U=(AX+BY)/R,\qquad W=(-BX+AY)/R.
\]
The equations become $U=C/R$ and $W^2=1-C^2/R^2=D/R^2$, which give \eqref{eq:intersectionpoints}. The angular derivative is $-A\sin\theta+B\cos\theta=BX-AY=-RW$, whose absolute value is $\sqrt D$. The Jacobian determinant of $X^2+Y^2-1$ and $AX+BY-C$ has absolute value $2\sqrt D$, expressing the same transversality condition.
\end{proof}
If $A=B=0$, the equation instead holds for every angle when $C=0$, and for no angle otherwise. For arbitrary surreal, rather than finite, angle parameters, each finite solution class of \eqref{eq:linecircle} repeats under the canonical period class $2\pi\Oz$.

\subsection{The algebra which does not disappear at tangency}
Put $d=D/R^2$. Over $\SC$, the coordinate algebra of the line--circle intersection is
\begin{equation}\label{eq:quadalgebra}
 \mathcal A_d=\SC[W]/(W^2-d).
\end{equation}
Concretely, its elements are the unique remainders $a+bW$, with $W^2=d$ defining multiplication; this also interprets the notation as a class algebra without taking a collection of proper-class cosets. It has basis $1,W$ and dimension two for every $d$, including $d=0$ and negative real $d$. When $d<0$ there are no real points, but the two complexified points still exist. When $d=0$ there is a single point with the nonreduced algebra of dual numbers. This is the precise sense in which the multiplicity survives the change in the number of real solutions.

\begin{theorem}[A collision-stable residue pairing]\label{thm:residuepairing}
Define $\Lambda_d:\mathcal A_d\to\SC$ by
\[
 \Lambda_d(a+bW)=b.
\]
The pairing $(f,g)\mapsto\Lambda_d(fg)$ is nondegenerate for every $d$, and in the basis $1,W$ its matrix is
\begin{equation}\label{eq:pairmatrix}
 \begin{pmatrix}0&1\\1&0\end{pmatrix}.
\end{equation}
If $d\ne0$ and $s^2=d$, then for every polynomial $g$,
\begin{equation}\label{eq:residuesum}
 \Lambda_d(g)=\frac{g(s)}{2s}+\frac{g(-s)}{-2s}
 =\frac{g(s)-g(-s)}{2s}.
\end{equation}
At $d=0$ the corresponding formula is $\Lambda_0(g)=g'(0)$. The trace of multiplication by $g$ is $\Lambda_d(2Wg)$.
\end{theorem}
\begin{proof}
The relation $W^2=d$ gives
\[
 \Lambda_d(1)=0,\quad\Lambda_d(W)=1,\quad\Lambda_d(W^2)=\Lambda_d(d)=0,
\]
proving \eqref{eq:pairmatrix}, whose determinant is $-1$. Reduction of $g$ modulo $W^2-d$ gives $g=a+bW$. Its values at $s,-s$ are $a+bs,a-bs$, so \eqref{eq:residuesum} equals $b$. These are exactly the two simple residues of $g(W)/(W^2-d)$. At $d=0$, reduction gives $g(0)+g'(0)W$, and the coefficient of $W^{-1}$ in $g(W)/W^2$ is $g'(0)$. Finally multiplication by $a+bW$ has matrix
\[
 \begin{pmatrix}a&bd\\b&a\end{pmatrix},
\]
whose trace is $2a=\Lambda_d(2W(a+bW))$.
\end{proof}
For infinitesimal positive $d$, the two residues in \eqref{eq:residuesum} may individually be infinite, while their sum remains well behaved. For example $g(W)=1$ has residues $\pm(2\sqrt d)^{-1}$, which cancel exactly. The function $g(W)=W$ gives two residues $1/2$, whose sum is $1$. This elementary model is the same structural phenomenon emphasized in the supplied finite-geometry and intersection papers: the quotient and its pairing are more stable than the list of separated simple points.

\subsection{A ramified trigonometric equation}
\begin{theorem}[The two infinitesimal roots at a cosine extremum]\label{thm:cosroots}
Let $\eta\in\mm_{\SC}$. If $\eta\ne0$, the solutions in the infinitesimal monad of zero of
\begin{equation}\label{eq:cosperturb}
 \Cos z=1-\eta
\end{equation}
are exactly
\begin{equation}\label{eq:cosrootformula}
 z_\pm=\pm2\arcsin_{\mathrm{inf}}\sqrt{\eta/2},
\end{equation}
where $\arcsin_{\mathrm{inf}}$ is the formal inverse-sine series and either square root may be selected. Both solutions are simple. If $\eta=0$, zero is the unique root in that monad, of multiplicity two. The expansion is
\begin{equation}\label{eq:cosrootexp}
 z_\pm=\pm\sqrt{2\eta}
 \left(1+\frac\eta{12}+\frac{3\eta^2}{160}
                  +\frac{5\eta^3}{896}+\cdots\right).
\end{equation}
For real $\eta>0$ the roots are real; for real $\eta<0$ they are purely imaginary.
\end{theorem}
\begin{proof}
The identity $1-\Cos z=2\Sin^2(z/2)$ holds in the infinitesimal monad. The formal series $\Sin(z/2)$ has nonzero linear coefficient and is a bijection from that monad to itself, with inverse $w\mapsto2\arcsin_{\mathrm{inf}}w$. All its formal inverse identities evaluate by the support lemma. Thus the equation is equivalent to $w^2=\eta/2$, giving precisely the two stated roots. At a nonzero root, $\Sin z=2\Sin(z/2)\Cos(z/2)\ne0$, so the derivative of \eqref{eq:cosperturb} is nonzero. At $\eta=0$, $\Cos z-1=-z^2/2+\cdots$ gives the double root. Expanding the inverse-sine series gives \eqref{eq:cosrootexp}; its coefficients are real, so a real positive square root produces real roots and a purely imaginary square root produces purely imaginary roots.
\end{proof}
This is an explicit one-variable preparation: after the analytic coordinate $w=\Sin(z/2)$, the equation is the quadratic $w^2-\eta/2=0$. It is not merely a first-order approximation to a root count. In particular the exponent $v(z_\pm)=v(\eta)/2$ is exact.

\subsection{A sharp inverse-cosine estimate}
\begin{theorem}[Conditioned inversion of cosine]\label{thm:conditioned}
Let $\theta\in(0,\pi)$, $c=\cos\theta$, and $s=\sin\theta>0$. Let $0\ne\varepsilon\in\No$ satisfy
\begin{equation}\label{eq:condition}
 \sigma:=v(\varepsilon)>2\kappa,
 \qquad \kappa:=v(s)\ge0.
\end{equation}
Then $c+\varepsilon\in(-1,1)$. If
$h=\arccos(c+\varepsilon)-\theta$, then
\begin{align}\label{eq:inversecos}
 h&=-\frac\varepsilon s-\frac{c\varepsilon^2}{2s^3}+\mathcal R,\\
 v(\mathcal R)&\ge3\sigma-5\kappa,\nonumber\\
 v(h)&=\sigma-\kappa,\qquad
 v(h+\varepsilon/s)\ge2\sigma-3\kappa.\nonumber
\end{align}
The strict threshold in \eqref{eq:condition} cannot be weakened to equality uniformly near an extremum.
\end{theorem}
\begin{proof}
The distance from $c$ to the nearer endpoint of $[-1,1]$ is
\[
 1-|c|=\frac{s^2}{1+|c|}.
\]
Condition \eqref{eq:condition} says that $\varepsilon/s^2$ is infinitesimal, hence its magnitude is smaller than this distance divided by $s^2$. Thus $c+\varepsilon\in(-1,1)$.

Set $u=\varepsilon/s^2$ and introduce $h=sq$. The addition formula gives
\begin{equation}\label{eq:normalizedinverse}
 \frac{\cos(\theta+sq)-c}{s^2}
 =c\frac{\cos(sq)-1}{s^2}-\frac{\sin(sq)}s
 =-q-\frac c2q^2+\frac{s^2}{6}q^3+\cdots.
\end{equation}
Every coefficient on the right is a polynomial with rational coefficients in the finite parameters $c,s^2$. Its linear coefficient is $-1$. Formal inversion therefore yields
\[
 q=-u-\frac c2u^2+u^3H(u),
\]
where every coefficient of $H$ is again a polynomial in $c,s^2$ with rational coefficients. Evaluation at infinitesimal $u$ is legitimate: write the finitely many parameters as ordinary standard parts plus infinitesimals; all expanded terms are supported in the positive monoid generated by those infinitesimals and $u$. Positive powers of $u$ ensure finite contribution at each exponent. Thus $H(u)$ is finite and $v(H(u))\ge0$.

The constructed $h=sq$ satisfies $h/s\in\mm_{\R}$. It therefore stays inside the same interior-angle branch: the distance to either endpoint is at least $s$, since $s=\sin\theta\le\min(\theta,\pi-\theta)$, with the endpoint comparison obtained from \cref{prop:smallangles}. It solves the desired equation, so strict monotonicity of cosine identifies it with the stated difference of inverse values. Multiplying the inverse series by $s$ gives
\[
 \mathcal R=s u^3H(u),\qquad v(\mathcal R)\ge\kappa+3(\sigma-2\kappa).
\]
The linear term has valuation $\sigma-\kappa$, while all later terms have strictly larger valuation. This proves every estimate in \eqref{eq:inversecos}.

For sharpness, let $\eta>0$ be infinitesimal and choose $c=1-\eta$. Then $s^2=2\eta-\eta^2$ has valuation $v(\eta)$. Taking $\varepsilon=\eta$ gives $\sigma=2\kappa$ and moves the target to $1$: the two real roots merge and the perturbed root is no longer simple. Taking $\varepsilon=2\eta$, with the same valuation, moves the target above $1$, so real solutions disappear. Consequently a uniform simple-root stability theorem cannot replace $>$ by $\ge$.
\end{proof}
The roles of the two powers of the small derivative are now explicit. One factor of $s^{-1}$ amplifies the displacement; the other measures how close the target is to a critical value. The leading displacement has order $\sigma-\kappa$, and the nonlinear correction has order at least $2\sigma-3\kappa$, matching the inverse-Jacobian pattern in \cite{InputIntersections}. Near a noncritical ordinary angle $\kappa=0$, the condition is simply that the target perturbation be infinitesimal.

\section{Finite trigonometric polynomials and exact Fourier identities}\label{sec:fourier}
Nothing in this section concerns an uncontrolled infinite Fourier series. The number of frequencies and sample points is finite, while the coefficients themselves may have arbitrarily complicated set-sized surreal supports.

\begin{theorem}[A sharp algebraic bound on trigonometric zeros]\label{thm:trigpoly}
Let $N\ge1$ be an ordinary integer and
\[
 P(\theta)=a_0+\sum_{n=1}^N\bigl(a_n\cos(n\theta)+b_n\sin(n\theta)\bigr),
 \qquad a_n,b_n\in\No,
\]
with at least one nonzero coefficient. In one finite phase circle $\OO_{\R}/2\pi\Z$, $P$ has at most $2N$ zeros, counted with their fine-analytic multiplicities. The bound is attained by $\sin(N\theta)$.
\end{theorem}
\begin{proof}
For $u=\cis\theta$, Euler's formula expresses $P$ as the Laurent polynomial
\[
 P=c_0+\sum_{n=1}^N(c_nu^n+c_{-n}u^{-n}),\qquad
 c_n=(a_n-ib_n)/2,\quad c_{-n}=(a_n+ib_n)/2.
\]
Multiplication by $u^N$ gives a nonzero polynomial $Q(u)$ of degree at most $2N$. Its roots on $\UU$ correspond bijectively to angle classes by \cref{thm:phase}. Locally $u'(\theta)=iu\ne0$, so the substitution preserves vanishing order; multiplication by $u^N$ also preserves it. A polynomial has total root multiplicity at most its degree. Finally $\sin(N\theta)$ has the $2N$ distinct simple classes $\theta=k\pi/N$, $0\le k<2N$.
\end{proof}
Infinitesimal coefficients can split a multiple zero into a cluster as in \cref{thm:cosroots}, but cannot violate the total-degree bound. On all of $\No$, canonical global trigonometric polynomials repeat their zeros with the period class $2\pi\Oz$, so a global count without a phase quotient would be a different and generally infinite assertion.

\begin{theorem}[Discrete Fourier inversion and Parseval over $\SC$]\label{thm:dft}
Let $M\ge1$ be an ordinary integer, $\zeta=\cis(2\pi/M)$, and $z_0,\ldots,z_{M-1}\in\SC$. Define
\[
 \widehat z_k=\frac1M\sum_{j=0}^{M-1}z_j\zeta^{-jk}.
\]
Then
\begin{equation}\label{eq:dft}
 z_j=\sum_{k=0}^{M-1}\widehat z_k\zeta^{jk},\qquad
 \sum_{j=0}^{M-1}|z_j|^2=M\sum_{k=0}^{M-1}|\widehat z_k|^2.
\end{equation}
\end{theorem}
\begin{proof}
The finite geometric-sum identity gives
\[
 \sum_{j=0}^{M-1}\zeta^{j\ell}
 =\begin{cases}M,&M\mid\ell,\\0,&M\nmid\ell.\end{cases}
\]
Substituting the definition of $\widehat z_k$ and using this identity proves inversion. Substitute the inversion formula and its conjugate into $\sum_j z_j\bar z_j$. Orthogonality removes every term with different frequencies and leaves $M\sum_k\widehat z_k\overline{\widehat z_k}$. Every sum is finite, and $\bar\zeta=\zeta^{-1}$, so the calculation is valid in $\SC$ without additional support or convergence conditions.
\end{proof}
All nonzero terms in the norm identity are positive surreals. In particular Fourier cancellation cannot destroy energy at a hidden infinitesimal order; the equality includes every Hahn coefficient after the finite products are formed.

For the real-valued trigonometric polynomial in \cref{thm:trigpoly}, the coefficientwise integral convention gives the parallel identity
\begin{equation}\label{eq:continuousparseval}
 \frac1{2\pi}\HInt_0^{2\pi}P(\theta)^2\,d\theta
 =a_0^2+\frac12\sum_{n=1}^N(a_n^2+b_n^2).
\end{equation}
To prove it, expand the finite square and apply ordinary sine--cosine orthogonality to each coefficient. Multiplication by fixed surreal constants is legitimate because Hahn products have finite convolution at each exponent. Equation \eqref{eq:continuousparseval} is an exact finite-frequency statement; it implies no completeness or convergence assertion for an arbitrary infinite Fourier expansion.

\section{Spherical and hyperbolic analogues}\label{sec:curved}
The same distinction between algebraic geometry and angle evaluation extends beyond the plane. We record two substantial analogues, proved from Gram identities. This section uses only three-dimensional vectors and finite matrix operations.

\subsection{Spherical sine and cosine laws}
Let $P,Q,R\in\No^3$ be unit vectors which are linearly independent. Define the minor central side angles
\[
 a=\arccos(Q\cdot R),\quad b=\arccos(R\cdot P),\quad c=\arccos(P\cdot Q).
\]
They belong to $(0,\pi)$. At $P$, the unit tangent vectors toward $Q$ and $R$ are
\[
 T_{PQ}=\frac{Q-\cos c\,P}{\sin c},\qquad
 T_{PR}=\frac{R-\cos b\,P}{\sin b}.
\]
Let $A\in(0,\pi)$ be their angle; define $B,C$ cyclically.

\begin{theorem}[Spherical trigonometry]\label{thm:spherical}
The spherical cosine and sine laws hold:
\begin{gather}\label{eq:sphericalcos}
 \cos a=\cos b\cos c+\sin b\sin c\cos A,\\
 \frac{\sin A}{\sin a}=\frac{\sin B}{\sin b}
 =\frac{\sin C}{\sin c}
 =\frac{|\det(P,Q,R)|}{\sin a\sin b\sin c}.\label{eq:sphericalsine}
\end{gather}
The cyclic cosine laws hold as well. On a sphere of arbitrary positive surreal radius $\rho$, the corresponding coherent minor arc lengths are $\rho a,\rho b,\rho c$.
\end{theorem}
\begin{proof}
The tangents are unit vectors perpendicular to $P$. Their dot product is
\[
 T_{PQ}\cdot T_{PR}
 =\frac{\cos a-\cos b\cos c}{\sin b\sin c},
\]
which proves \eqref{eq:sphericalcos}. In the tangent plane, the magnitude of the determinant of two unit tangent vectors is their angle sine. As $P$ is a unit normal,
\[
 \sin A=|\det(P,T_{PQ},T_{PR})|
 =\frac{|\det(P,Q,R)|}{\sin b\sin c}.
\]
The usual cross-product identity used here is a polynomial consequence of the Euclidean Gram determinant and is valid over $\No$. Dividing by $\sin a$ and relabeling proves \eqref{eq:sphericalsine}. Scaling the circle containing each great-circle arc gives the arc-length statement from \cref{prop:arc}.
\end{proof}
Linear independence excludes great-circle degeneracy and ensures that the angle sines and side sines in the denominators are nonzero. The theorem still allows all three points to be infinitesimally close, so it includes spherical triangles on arbitrarily small angular scales.

\subsection{Hyperbolic trigonometry with infinite distances}
Use the Lorentz form
\[
 L(P,Q)=P_0Q_0-P_1Q_1-P_2Q_2
\]
and the upper unit hyperboloid
\[
 \mathbb H^2_{\No}=\{P\in\No^3:L(P,P)=1,\ P_0>0\}.
\]
For $x\ge1$, define
\[
 \arcosh x=\log_{\No}(x+\sqrt{x^2-1}).
\]
The real exponential identities show that this is the inverse of $\cosh$ on $\No_{\ge0}$. For $P,Q$ on the hyperboloid, $L(P,Q)\ge1$, with equality only when $P=Q$. Indeed, writing the spatial Euclidean norms as $p,q$, Cauchy--Schwarz gives
\[
 L(P,Q)\ge\sqrt{1+p^2}\sqrt{1+q^2}-pq\ge1,
\]
where the second inequality follows by squaring and using
$(1+p^2)(1+q^2)-(1+pq)^2=(p-q)^2$. Equality also requires equality in spatial Cauchy--Schwarz, which gives the same point.

For linearly independent $P,Q,R\in\mathbb H^2_{\No}$, set
\[
 a=\arcosh L(Q,R),\quad b=\arcosh L(R,P),\quad c=\arcosh L(P,Q).
\]
These positive distances may be infinite. On the tangent space $L(P,V)=0$, the form $-L$ is positive definite. For if $P=(P_0,p)$ and $V=(V_0,v)$, then $V_0=p\cdot v/P_0$, and
\[
 -L(V,V)=|v|^2-(p\cdot v)^2/P_0^2\ge|v|^2/P_0^2>0
\]
for nonzero $V$. Define the angle $A\in(0,\pi)$ using this positive metric between
\[
 T_{PQ}=\frac{Q-\cosh c\,P}{\sinh c},\qquad
 T_{PR}=\frac{R-\cosh b\,P}{\sinh b},
\]
and define $B,C$ cyclically.

\begin{theorem}[Hyperbolic sine and cosine laws]\label{thm:hyperbolic}
One has
\begin{gather}\label{eq:hyperboliccos}
 \cosh a=\cosh b\cosh c-\sinh b\sinh c\cos A,\\
 \frac{\sin A}{\sinh a}=\frac{\sin B}{\sinh b}
 =\frac{\sin C}{\sinh c}
 =\frac{|\det(P,Q,R)|}{\sinh a\sinh b\sinh c}.\label{eq:hyperbolicsine}
\end{gather}
\end{theorem}
\begin{proof}
The tangent vectors have $L$-norm $-1$. Their angle cosine is
\[
 -L(T_{PQ},T_{PR})
 =\frac{\cosh b\cosh c-\cosh a}{\sinh b\sinh c},
\]
which proves the first law. For the second, the Lorentz Gram matrix of the columns $P,T_{PQ},T_{PR}$ is
\[
 \begin{pmatrix}1&0&0\\0&-1&-\cos A\\0&-\cos A&-1\end{pmatrix}.
\]
Its determinant is $\sin^2 A$. Since the matrix of $L$ has determinant $1$, the same Gram determinant is the square of $\det(P,T_{PQ},T_{PR})$. Taking positive square roots and substituting the tangent vectors yields
\[
 \sin A=\frac{|\det(P,Q,R)|}{\sinh b\sinh c}.
\]
Cyclic relabeling proves \eqref{eq:hyperbolicsine}.
\end{proof}
Unlike circular angle functions, the hyperbolic functions use the ordered real surreal exponential without discarding a purely infinite part. Infinite distances therefore represent genuine large hyperbolic separation; the angles between tangent directions remain finite. We have proved the displayed metric identities, not a general theory of hyperbolic area or of all possible notions of non-Archimedean geodesic completeness.

\section{Boundaries of the theory and failure of naive extensions}\label{sec:limits}
\subsection{Hahn sums are not fine-topological limits}
The fine topology uses all balls $|z-a|<r$ with positive surreal $r$. For every set $S\subseteq\No_{>0}$ there is a positive surreal smaller than every member of $S$, for example the cut $\{0\mid S\}$ when $S$ is nonempty. This makes the full class topology much finer than any topology controlled by an ordinary sequence of scales.

\begin{proposition}[Set-sized convergence degeneracy]\label{prop:topology}
Every set-sized subset of $\SC$ is closed and discrete in the fine topology. A net with a set-sized directed index set which converges in that topology is eventually equal to its limit. In particular a nonconstant ordinary real-parameter path cannot be fine-continuous on a connected real interval.
\end{proposition}
\begin{proof}
For a set $A\subseteq\SC$ and a point $p$, take a positive surreal smaller than all nonzero distances $|p-a|$, $a\in A$. Its ball contains no point of $A$ other than possibly $p$, proving both closedness and discreteness. For a convergent net use a positive lower bound for all its nonzero distances from the alleged limit; eventually entering that ball means eventual equality. The image of an ordinary real interval is a set, hence a discrete subspace. A continuous image of a connected space is connected, so it must be a singleton.
\end{proof}
For example the strong series $\sum_{n\ge0}\omega^{-n}$ exists, but its finite partial sums do not converge to it in this topology. Their nonzero tails have a common smaller positive surreal bound. This is why all summation arguments in the article check supports and finite contribution, and why \cref{prop:arc} explicitly defines a coefficientwise functional.

\subsection{A global Maclaurin series is not the definition}
At $x=\omega$, the terms of the putative sine series
\[
 \sum_{n\ge0}\frac{(-1)^n\omega^{2n+1}}{(2n+1)!}
\]
have valuations $-(2n+1)$, a strictly decreasing sequence. Their support union is not well ordered, so this is not a strong sum. The canonical value $\Sin\omega=0$ is supplied by the integer-part normalization, not by an illicit evaluation or rearrangement of this series. The same warning applies at infinite complex arguments to the exponential series.

Local analyticity does not repair the missing global uniqueness. The characters in \cref{thm:characters} give different exponentials with the same ordinary restriction and the same differential equation. Conversely, the existence of these alternatives does not negate the canonical status of the initial-integer-part construction: canonicity refers to that additional structure, not to uniqueness under the weaker analytic hypotheses.

\subsection{Why demanding coherence at every scale is too strong}
A Hahn-coherent analytic datum on an ordinary complex domain $U$ is a family
\[
 F=\sum_{\lambda\in S}f_\lambda t^\lambda,
 \qquad f_\lambda\in\mathcal O(U),
\]
with one well-ordered support $S$. It evaluates at finite $c+\eta$ by lifting each coefficient and summing. This is the notion in \cite{InputAnalysis}; it is stronger than independently prescribing local power series on different monads.

Suppose that a global function $f$ had, at every positive surreal radius $r$, such a chart for $w\mapsto f(rw)$ on the infinitesimal thickening of a fixed ordinary neighborhood of zero. Write $a_n=f^{(n)}(0)/n!$. The support of every $a_nr^n$ would then be contained in the chart's one well-ordered support, because ordinary differentiation introduces no new Hahn exponents.

If infinitely many $a_n$ are nonzero, choose a positive surreal $\Delta$ larger than four times every $|v(a_n)|$, and take $r=t^{-\Delta}$. For successive nonzero indices $m>n$,
\[
 v(a_mr^m)-v(a_nr^n)
 =v(a_m)-v(a_n)-(m-n)\Delta<0.
\]
This produces an infinite decreasing sequence inside the alleged chart support, a contradiction. Thus every such all-scale coherent function has only finitely many nonzero Taylor coefficients at zero. In particular neither the global exponential nor global sine can satisfy this condition. The stronger polynomial classification, including global propagation of that finite germ, is proved in the supplied analysis paper; the displayed obstruction already suffices here.

This is a restriction on uniform analytic descriptions, not a failure of the algebraic trigonometric identities or of the arbitrary-radius circle geometry. A circle with huge radius uses the map $R\cis\theta$ with \emph{finite angle}; it does not require the function $w\mapsto\Sin(Rw)$ to possess one Hahn support over an entire ordinary angular domain.

\subsection{What has and has not been generalized}
The geometric identities proved above are equalities and inequalities of actual surreal numbers. They include the infinitesimal information suppressed by standard parts, and they allow arbitrary positive length scales. The local estimates use finite-angle Taylor lifting; the global function theory uses a specified, published integer-part normalization. These are compatible but logically distinct achievements.

No arbitrary infinite polygon, full-class Riemann integral, ordinary sequential Fourier convergence theorem, or uniqueness theorem for all fine-analytic differential equations has been asserted. The spherical and hyperbolic sections give metric identities under explicit nondegeneracy hypotheses, not a complete global geometric or measure-theoretic foundation. The algebraic collision model is fully computed, but is not a proof of the general multivariate intersection results in the supplied papers. Nor is a solution of the separate robustness questions about all possible initial fields claimed.

\section{Conclusion}
The unit circle is the bridge between surcomplex analysis and geometry. Its exact decomposition into an ordinary direction and an infinitesimal angular correction proves that every direction has a finite angle. Dot-product and determinant identities then recover the classical triangle and circle laws at infinitesimal and infinite scales. The same phase map extends canonically to every surreal argument after choosing the initial omnific integer part, yielding the established global surcomplex exponential and its noncyclic period class.

The extra geometric content lies in retaining the infinitesimal orders. Phase is a local valuation isometry, but displacement also includes the radial scale. A nearly flat triangle has a square-root angular defect and reciprocal square-root circumradius, governed by a relative side-gap threshold. At a trigonometric tangency, the two roots merge without destroying the dimension or residue pairing of the intersection algebra. Away from the exact collision, the sharp inverse-cosine estimate explains precisely how a small derivative amplifies target perturbations. These results demonstrate a useful trigonometry that is neither a formal repetition of real formulas nor dependent on an unjustified global convergence theory.

\appendix
\section{Notation and domain guide}\label{app:notation}
\begin{center}\small
\begin{longtable}{@{}L{0.24\textwidth}L{0.69\textwidth}@{}}
\toprule
Notation & Meaning and domain\\
\midrule
\endhead
$\No$, $\SC=\No[i]$ & The surreal real class field and its complexification.\\[0.4em]
$t^\gamma=\omega^{-\gamma}$ & Conway monomials, with increasing Hahn exponent convention.\\[0.4em]
$v(z)$, $\lc(z)$ & Least support exponent and leading complex coefficient of $z\ne0$; $v(0)=+\infty$.\\[0.4em]
$\OO_{\R}$, $\mm_{\R}$ & Finite real surreals and real infinitesimals; analogous subscripts $\SC$ for the complexification.\\[0.4em]
$\st(z)$ & Ordinary real or complex standard part of a finite element.\\[0.4em]
$\Pi$, $\fin(x)$ & Purely infinite normal forms; additive finite-part projection retaining the infinitesimal part.\\[0.4em]
$\Oz=\Pi\oplus\Z$ & The omnific integer part. Its finite elements are precisely $\Z$.\\[0.4em]
$\UU$ & The algebraic unit circle $\{u:u\bar u=1\}$; every point has finite coordinates.\\[0.4em]
$\sin,\cos,\cis$ & Taylor-lifted functions at finite arguments; $\cis\theta=\cos\theta+i\sin\theta$ for finite real $\theta$.\\[0.4em]
$\Sin,\Cos,\Exp$ & Canonical global functions on $\SC$; agree with the finite lifts on finite arguments.\\[0.4em]
$\exp_{\No},\log_{\No}$ & The ordered real surreal exponential and logarithm.\\[0.4em]
$\arccos,\arcsin,\arctan$ & Real inverse branches with ranges $[0,\pi]$, $[-\pi/2,\pi/2]$, and $(-\pi/2,\pi/2)$; arctangent's domain is all $\No$.\\[0.4em]
$a\prec b$, $a\sim b$ & $a/b$ is infinitesimal; respectively $a/b-1$ is infinitesimal.\\[0.4em]
$\langle u,v\rangle$, $[u,v]$ & Real dot product and oriented determinant in $\SC\cong\No^2$.\\[0.4em]
$a,b,c;\alpha,\beta,\gamma$ & Triangle sides and opposite interior angles; angles are finite even when sides are not.\\[0.4em]
$\Delta,R,r,s$ & Triangle area, circumradius, inradius, and semiperimeter; $s$ is also explicitly reused as a local sine or square root in \S\ref{sec:collisions}.\\[0.4em]
$\HInt$ & Coefficientwise Hahn integral, with analytic endpoint corrections where specified; not an ordinary fine-topological Riemann integral.\\
\bottomrule
\end{longtable}
\end{center}

\section{Support audit and logical dependencies}\label{app:support}
Every analytic use of an infinite series in the article reduces to one of the following cases.
\begin{center}\small
\begin{tabular}{@{}L{0.36\textwidth}L{0.57\textwidth}@{}}
\toprule
Construction & Reason the sum is legitimate\\
\midrule
Finite sine, cosine, exponential & Ordinary values are evaluated first; positive-support infinitesimal powers are then summed.\\[0.5em]
Unit-circle logarithm & Its argument belongs to $1+\mm_{\SC}$, so the logarithm series has positive-support control.\\[0.5em]
Arctangent at infinite slope & Its reciprocal argument is infinitesimal; the expansion is at zero in that reciprocal.\\[0.5em]
Square-root and inverse-sine expansions & The small parameter or its square root is infinitesimal; all supports lie in the corresponding positive monoid.\\[0.5em]
Conditioned inverse cosine & The normalized parameter $\varepsilon/s^2$ is infinitesimal, and coefficients are polynomials in finitely many finite parameters.\\[0.5em]
Coefficientwise circular integrals & The pulled-back speed and sector integrand are constants; more general finite trigonometric sums preserve a common Hahn support.\\[0.5em]
Global phase extension & No new infinite sum at an infinite angle is taken. Its normal form is split, and a finite-angle value is selected.\\
\bottomrule
\end{tabular}
\end{center}
Finite algebraic identities, determinants, Fourier transforms, and finite-degree polynomial root counts require no additional summability hypothesis. Square roots use real closedness for positive real quantities and algebraic closedness for complex quantities. The class-wide topology arguments use the availability of a scale smaller than every member of a set; they should not be silently transferred to a single fixed Hahn field with its own valuation topology.

The foundational dependencies can be summarized without circularity. Normal forms and the support lemma give Taylor lifts and infinitesimal logarithms. These give the phase theorem and the finite inverse branches. Dot products and the phase theorem give triangle angles and the angle sum; the coordinate calculation gives the sine law. Taylor leading terms applied to those exact identities yield the valuation and degeneration theorems. The sharp inverse theorem uses only finite trigonometry and a normalized formal inverse. The canonical global extension is separate: it adds the real surreal exponential and the omnific integer-part normalization, but is not needed to justify finite triangle angles.

\section{Verification notes and reproducibility}\label{app:verification}
The companion script \texttt{verify\_examples.py} performs exact symbolic checks with SymPy. It checks finite Taylor jets of the phase identities, rational stereographic parameterization, the Gram and Ptolemy polynomial identities, Heron's factorization, the line--circle formulas, the collision algebra, the inverse-cosine coefficients, the fractional-power examples, and an exact discrete Fourier instance. The generated \texttt{verification\_report.txt} records the checks actually executed; all 43 checks passed in the build accompanying this article.

These checks detect sign, factor, and coefficient errors, but do not prove the support lemma, general root multiplicities, or the sharp valuation estimates. Those rely on the mathematical proofs. No proof-assistant verification is claimed. The standalone LaTeX source includes the bibliography; build instructions and tested software versions are in the README.

\begin{thebibliography}{99}
\addcontentsline{toc}{section}{References}

\bibitem{Conway}
John H. Conway,
\emph{On Numbers and Games},
Academic Press, 1976; second edition, A K Peters, 2001.
Normal forms, the ordered-field structure, and omnific integers.

\bibitem{Gonshor}
Harry Gonshor,
\emph{An Introduction to the Theory of Surreal Numbers},
London Mathematical Society Lecture Note Series \textbf{110},
Cambridge University Press, 1986.
Surreal normal forms, real closedness, and the real exponential.

\bibitem{Neumann}
B. H. Neumann,
``On ordered division rings,''
\emph{Transactions of the American Mathematical Society} \textbf{66} (1949), 202--252.
The ordered-support summability lemma.

\bibitem{Alling}
Norman L. Alling,
\emph{Foundations of Analysis over Surreal Number Fields},
North-Holland Mathematics Studies \textbf{141},
North-Holland, 1987.
Formal analytic evaluation and local analysis over surreal fields.

\bibitem{vdDE}
Lou van den Dries and Philip Ehrlich,
``Fields of surreal numbers and exponentiation,''
\emph{Fundamenta Mathematicae} \textbf{167} (2001), 173--188.
\href{https://doi.org/10.4064/fm167-2-3}{doi:10.4064/fm167-2-3}.
Erratum: \textbf{168} (2001), 295--297.
The present article uses restricted analytic evaluation and exponential-field properties, not the ordinal bounds addressed by the erratum.

\bibitem{EK}
Philip Ehrlich and Elliot Kaplan,
``Surreal ordered exponential fields,''
\emph{The Journal of Symbolic Logic} \textbf{86}, no. 3 (2021), 1066--1115.
\href{https://doi.org/10.1017/jsl.2021.59}{doi:10.1017/jsl.2021.59};
\href{https://arxiv.org/abs/2002.07739}{arXiv:2002.07739}.
In particular \S11: canonical trigonometric functions from an initial integer part and canonical surcomplex exponentiation. The arXiv version consulted is v3, June 22, 2021. The paper credits an earlier idea of using $\Oz$ for trigonometric functions to Martin Kruskal.

\bibitem{InputAnalysis}
\emph{Surcomplex Analysis: Infinitesimal Calculus, Hahn-Coherent Holomorphy, and Contour Theory}.
Supplied manuscript, September 21, 2026;
file \texttt{surcomplex\_analysis(2).tex}.
Used for its stated distinctions among fine-local analyticity, strong summability, scalar lifts, coherent contours, and global phase choices. This is user-supplied material, not identified here as a peer-reviewed publication.

\bibitem{InputGeometry}
\emph{Finite Geometry of Hahn-Coherent Surcomplex Zeros:
Conservation of Multiplicity, Normal Forms, and Collision-Stable Residues}.
Supplied manuscript, September 21, 2026, in
\texttt{surcomplex\_finite\_geometry.zip}.
Motivates retaining a finite intersection algebra and its residue pairing through a root collision; the quadratic instance here is proved directly.

\bibitem{InputIntersections}
\emph{Finite Intersections in Surcomplex Analysis:
Support-Controlled Division, Residue Duality, and Sharp Root Stability}.
Supplied manuscript, September 21, 2026, in
\texttt{surcomplex\_intersections.zip}.
Its \S9 treats conditioned root stability with the threshold $\sigma>2\kappa$. The inverse-cosine specialization and second-order expansion in this article have independent proofs.

\end{thebibliography}
\end{document}
